\chapter[Against Einstein's Relativity: A Doxographical Analysis]{Against Einstein's Relativity:\\A Doxographical Analysis\label{chap_AER}}
\epigraph{I cannot believe---and I say this with all the emphasis of which I am capable---that there can ever be any good excuse for refusing to face the evidence in favour of something unwelcome. It is not by delusion, however exalted, that mankind can prosper, but only by unswerving courage in the pursuit of truth.}{---Bertrand Russell, \textit{Fact and Fiction}}
\section{Introductory Remarks\label{sec_IR}}
The primary goal of this investigation, maintained from the outset, has been to find a natural description of reality that is implicitly consistent with our observations---not by trying to rectify any apparently paradoxical or conflicting facts within an assumed framework, but by attempting to identify and analyse together all the basic facts that we can know objectively, and using them as indicators of what would therefore seem, prior to forming any prejudice, even according to any logical reasoning, to be the most natural interpretation; i.e., not to allow only for na\"{i}vely objective deduction of only certain specific facts, neglecting that the resulting theory might ultimately be at odds with certain others, and then opting to ignore such problems, but also to base any subsequent deduction on some common sense and basic observational facts, which are relevant to our description of reality. 

The idea is, that in specific instances regarding our observations of reality, multiple interpretations can be made logically,---some of which do not however fit naturally into the greater picture, as inconsistencies are later discovered. Logically, therefore, in the grand scheme, such interpretations are therefore found not to fit into a natural description of reality. It might be thought, that if the whole picture could be considered objectively, such paradoxical interpretations would never enter into our theories, as they would be immediately removed by such logic. We are thus led to consider Laplace's demon \cite{Laplace1951}, who would know, e.g., which interpretation of quantum theory is actually correct, by considering and comprehending all of reality at once. Because it is beyond human capability to do this, we are limited to making logical interpretations of more specific observations.

But this is inherently problematic; for if we would, e.g., regard specific observations always with complete independence, then we should never be able to decipher which, of the many logical interpretations that would be allowed, would actually be the correct one; and even if we lessen the extremeness of this requirement, so that different objective measurements would be considered together in our deductions, we can never consider \textit{everything}; so we cannot determine Truth---if such exists, which we take \textit{in principle} to be so---purely through objective deduction in considering the logical interrelatedness of multiple observed facts. 

On the other hand, we have our sensory perception of Nature, which, although a fact of observation, is difficult to rationally inspect; in particular, it is often difficult to distinguish the conclusions that are drawn intuitively from such perceptions---and, indeed, whether those are not perceived at all, but conceived only by the mind---from the basic facts of perception. Regardless, we must lay some trust in our own basic sensory perceptions, and, if we have seen that they are common ones, suppose that they may result from some property of Nature we are thus observing, and allow them to guide our interpretations of reality while keeping the limitations of human perception and understanding in mind; for the facts contributing to common sense seem to be critically important, and may be unobservable otherwise. 

Therefore, we say that when a scientist makes some observation, which could be logically interpreted many different ways, they are not being properly objective and truthful if they would say, `The truth could as well be any of these things, which I therefore shan't commit to any further,' if they know that certain of these interpretations are only allowed because the observation is being considered singularly, or even in conjunction with some others, but in denial of some aspect of common sense, merely brushing it aside; for if they would be truly objective, they would consider together all relevant observations, including basic sensory perception, and make their interpretation based on whichever thus seems most logical; for, along with the many other principles to which one must continually remain consistent, in order to be truly objective one must also be true to oneself.

It must be noted that any such examination of beliefs and biases, which aim is to discover their Natural foundations and include them in physical theory, is very different from yielding to common prejudice in theoretical development---whether that prejudice initially came as the result of critical thinking or otherwise; for only a theory that is consistent with \textit{all} empirical data is worthy of any amount of such prejudice, and then only so long as no subsequent inconsistencies are discerned, and only to the point that its basic principles can be accepted as the True facts of Nature.

The problem on which the next two chapters are centred, is that of the True metaphysical nature of Time, and how that is described in relativistic cosmology. For the true nature of Time has long been a topic of study in natural philosophy, and has by no means been resolved by relativity theory---the standard interpretation of which---if any one standard can really be said to exist---seems to lead to more paradoxes and inconsistencies than resolutions; therefore, in order to arrive at a general relativistic theory that could be used to describe the large-scale structure and evolution of our Universe, an appropriate view of the nature of time, and how that is described by the standard spacetime formalism of general relativity theory, must first be developed.

When this has been accomplished, the theory will eventually resolve as one consistent with the mathematical formalism of relativity theory, but in which the many well-known inconsistencies and problems, the apparent paradoxes which have commonly led to \textit{ad hoc} conjectures, etc., identically vanish. It is, by the way, the only theory of time, by which the dynamical cosmology described in Chapter~\ref{chap_ANC} makes any sense; but also the one by which the structure, evolution, and \textit{origin} of the Universe has potential to be truly explained \textit{a priori}, rather than given, as it is with the standard cosmology, as a description \textit{a posteriori}, which, although more general in application, is no more than that, and, in fact, which is the result of an interpretation of relativity theory requiring that the Universe does not actually \textit{evolve} at all, but is laid out, once and for all, as a singular four-dimensional continuum of events, which does not admit any \textit{real} dynamics.

The contrary picture I will therefore describe---which must be described, because it is the only one through which any meaningful \textit{understanding} can come of the foundational principles of relativistic cosmology---which, when stated below, will correctly be understood to be only superficially and trivially satisfied by the standard theory,---is not altogether new. In fact, it is closely related to the cosmological theory that Weyl developed in the early-1920s, over which he did battle with Einstein \cite{Weyl1924}, and also to a theory that was eventually reasoned out more than a millennium earlier, by St.~Augustine---which is really only extended here, by requiring a \textit{physical} fourth dimension of reality in that theory, and adding to it the formal mathematical framework of relativity theory. 

Although an argument for that theory must ultimately be given, the greater goal of this analysis is that it should provide a clear \textit{explanation} of this one theory. But in order to give a clear account of the meaning of this theory, and thus explain the general idea of what it is, which ultimately requires knowledge of what it is not, the discussion cannot be a complete logical deduction of all possibilities, which would ultimately obscure the straight and narrow path that is necessary in order to meet that specific purpose. Therefore, although an objective deduction of all the facts that are known, or may be immediately inferred, as well as the various philosophical theories that have been proposed throughout history, must be possible, the argument that is so vital to the explanation, is presented by means of whatever method of logical inquiry seems to be the most direct means to that end, at each point, so that we are able to maintain logical consistency while addressing only those points which seem to be the most pertinent to the problem.

In order to properly present such an argument, it is important to start at the beginning; for without a proper account of the rudiments of our modern theory, only a shallow explanation is possible; as any one wishing to fully understand a theory must know the reasoning behind it. Furthermore, it is at once obvious that in analysing the very roots of discovery, objective hindsight and modern intelligence should serve to dispel any errors that might have entered into a theory through contemporary prejudice, or to clearly see at which points subtle choices were actually made as such, when only one clear path may have been originally noticed---or, in general, to navigate the way along a path which was not originally taken, but which now seems to be more fruitful, as the most conducive to the discovery and formation of a better understanding of Nature. 

For in these regards we find ourselves in common sentiment with one of the greatest physicists in history, who wrote that `It is of great advantage to the student of any subject to read the original memoirs on that subject, for science is always most completely assimilated when it is in the nascent state' \cite{Maxwell1892}. For it is not from the epitomes of the original arguments, but from the prior reasonings through which those conclusions were drawn, that we stand to gain the most insightful perspective from which to improve on our understanding. In western philosophy, it appears that the first meaningful speculations about the Nature of Time came from opposing viewpoints, described by the Eleatics, who followed Parmenides, and by Heraclitus, who both flourished just prior to the turn of the fifth century \textsc{b.c.}, and whose works now survive only in fragmentary quotations and secondary accounts, as with all of the pre-Socratic philosophers. 

Before addressing the relevant aspects of these two theories, which we'll do in \S\S~\ref{sec_OH} -- \ref{sec_OE}, I think it should help to clarify the greater scope of this chapter and the next---and further, to see precisely how important to modern theory the works of Heraclitus and the Eleatics really are---to contrast it with a statement to which it is intimately linked, made by Sextus Empiricus, whose collective works, written in the late second century \textsc{a.d.}, are an important source of pre-Socratic philosophy \cite{Sextus2000} (III.65):
\begin{quotation}
The most fundamental positions on motion have, I think, been three in number. Common sense and some of the philosophers suppose that there is such a thing as motion; Parmenides and Melissus and some others think that there is not; and the Sceptics have said that there no more is than is not---so far as appearances go there seems to be motion, so far as philosophical argument goes it is unreal.
\end{quotation}
Sextus then sets out to show that the stance of the Sceptics is the right one.

In contrast, I shall maintain throughout, that the Sceptics' viewpoint is perfectly hypocritical. I say `maintain', because I will not attack this view directly, but will carefully analyse the other two theories, as the only two formal possibilities, and eventually discuss their relevance to general relativity theory. And I say `hypocritical' because the Sceptical stance, as it is here described, is \textit{essentially} Eleatic, but it pretends not to be the same as Eleatism, based on a concession that motion must be apparent, even though this was also recognised by the Eleatics. In principle, one must admit that a philosophical theory cannot \textit{be} one way and \textit{appear} another way, without still \textit{being} what it essentially \textit{is}.

So it is, that Sextus has really only recognised two different essential positions. If any meaningful distinction could be said to exist between the Eleatic and Sceptic views, as he has tried to say, it must be only, that on this point Scepticism regresses from the sophistication of its predecessor, and forgets what it is \textit{in principle}, which na\"{i}vet\'{e} inevitably gives rise to the acceptance of apparent paradoxes as genuine, yet abstract, truths---e.g., `that there no more is than is not' motion.

This can be related to more recent philosophic and scientific theories, beginning with Immanuel Kant's transcendental, or critical idealism: viz., as Kant correctly realised, nothing of reality can be observed, except its \textit{ideal}\footnote{This is meant in the adjectival sense of \textit{idea}, as in the French \textit{id\'{e}el}, or German \textit{ideell}.} aspects; so he refused to consider any other possible aspect of reality, other than that which may be subjected to science. Kant did not deny the existence of a metaphysical reality which would give rise to the observed \textit{ideality}, but, since no scientific description can be ultimately dissociated from an underlying metaphysical world-view, in denying the importance of describing and understanding the basic metaphysical noumena by pure reason, and contending that truth is only in experience, his theory becomes effectively equivalent to the philosophies of later idealists, such as Francis Herbert Bradley, John Ellis McTaggart, and Einstein, who ultimately committed to the metaphysical view that the ideality of experience is a four-dimensional block Reality.

%But, since the formal description of relativity theory, according to its original interpretation, \textit{is} that of a strictly determined four-dimensional ideality, relativists have often followed the Sceptical tradition; for in many cases, e.g.~in geometrodynamics, they reject the view of their predecessors, that time is purely an illusion, but retain the basic metaphysical framework, or formal interpretation of the mathematical theory. The essentially paradoxical nature of the current conception of relativity theory is in fact fully embodied in the partial completeness of geometrodynamics, and is manifest in the introductory section of the seminal paper by Richard Arnowitt, Stanley Deser, and Charles Misner \cite{Witten1962,Arnowitt2008}, as well as in the standard block spacetime description of dynamic gravitational collapse, which will be discussed in \S~\ref{sec_TBNP}. As another example, one can compare the theory of a causal structure, set out in \cite{Penrose1967,Carter1971}, which is formally based on the original interpretation of spacetime in relativity theory, with the dynamical description of reality given by Stephen Hawking, throughout \textit{A Brief History of Time} \cite{Hawking1988}; or, one might compare what is clearly Hawking's motive in \cite{Hawking1969}, in which he proves that a certain class of spacetime manifolds must admit a cosmic time function---which is essentially a Heraclitean concept, if for no reason other than the origin of that concept in the cosmological theory with Weyl---with the motive behind the theory of causal structures, which is to define what is necessary for an Eleatic spacetime manifold to \textit{appear} Heraclitean to any observer, for all `time'.

For, in Einstein's case, he did recognise the fact, that according to the interpretation of the relativity of simultaneity in which the simultaneous realities of two relatively moving observers are supposed to exist concurrently, as formally distinct sets of nevertheless real events, not only do those particular two realities coexist, but so too must the entire spacetime continuum of events, as a strictly determined block, which is all that may be investigated by science. As Louis de~Broglie once put it \cite{Einstein1951},
\begin{quotation}
In space-time, everything which for each of us constitutes the past, the present, and the future is given in block, and the entire collection of events, successive for us, which form the existence of a material particle is represented by a line, the world-line of the particle. Moreover, this new conception defers to the principle of causality and in no way prejudices the determinism of phenomena. Each observer, as his time passes, discovers, so to speak, new slices of space-time which appear to him as successive aspects of the material world, though in reality the ensemble of events constituting space-time exist prior to his knowledge of them. Although overturning a large number of the notions held by classical physics, the special theory of relativity may in one sense be considered as the crown or culmination of that physics, for it maintains the possibility of each observer localizing and describing all the phenomena in the schema of space and time, as well as maintaining the rigorous determinism of those phenomena, from which it follows that the aggregate of past, present, and future phenomena are in some sense given \textit{a priori}. \ldots
\end{quotation}
As discussed in \S~\ref{sec_OIS}, this was well-known to Einstein, and a formal argument has been given in special relativity theory by Wim Rietdijk \cite{Rietdijk1966} and Hilary Putnam \cite{Putnam1967}. For a more general discussion, one may refer to Kurt G\"{o}del's paper in \cite{Einstein1951} as well, which had previously examined the problem in relation to more general cosmological models. Some more recent investigations can be found, e.g., in \cite{Petkov2007}.

Of course, not everyone has understood that, aside from the `arrow of time' which provides its causal structure, the standard spacetime formalism should describe nothing more than points in a four-dimensional block. Therefore, by often advancing forms of subjective idealism that mirror the philosophy of George Berkeley, other descriptions of the theory seem to have regressed from the scientific idealism of Kant, formalised by the mathematical theory, which actually requires strict determinism after the standard interpretation of the relativity of simultaneity has been assumed; e.g., this is the case with the descriptions given by Eddington in many of his works,---as in \textit{The Nature of the Physical World} \cite{Eddington1929}, where he provides a paltry description of the meaning of relativity theory involving a subjectively-defined Here-Now and its corresponding light-cone, along with an imprecise neutral wedge, which he called Absolute Elsewhere---which he illustrated with the following hypothetical example:
\begin{quote}
Suppose that you are in love with a lady on Neptune and that she returns the sentiment. It will be some consolation for the melancholy separation if you can say to yourself at some---possibly prearranged---moment: `She is thinking of me now.' Unfortunately a difficulty has arisen because we have had to abolish Now. There is no absolute Now, but only the various relative Nows differing according to the reckoning of different observers and covering the whole neutral wedge which at the distance of Neptune is about eight hours thick. She will have to think of you continuously for eight hours on end in order to circumvent the ambiguity of `Now.'
\end{quote}
An example which illustrates the problem with such a subjective idealist scenario is the `Andromeda paradox', stated by Penrose in \textit{The Emperor's New Mind} \cite{Penrose1989}, according to which a group of Andromedeans might, in one frame, be considering whether to invade the Earth, with the ultimate decision being yet undetermined, while, in some relatively moving frame, they could already be on their way. 

But even some descriptions which don't enter into discussion about the paradoxical consequences that come with thinking of spacetime as anything less than a solidly determined block, often do not say explicitly that this must be so. For example, in \textit{The ABC of Relativity} Bertrand Russell came the \textit{closest} to saying this only at the end of his chapter on spacetime, where he wrote \cite{Russell1969},
\begin{quotation}
We may now recapitulate the reasons which have made it necessary to substitute `space-time' for space and time. The old separation of space and time rested upon the belief that there was no ambiguity in saying that two events in distant places happened at the same time; consequently it was thought that we could describe the topography of the universe at a given instant in purely spatial terms. But now that simultaneity has become relative to a particular observer, this is no longer possible. What is, for one observer, a description of the state of the world at a given instant, is, for another observer, a series of events at various different times, whose relations are not merely spatial, but also temporal. For the same reason, we are concerned with \textit{events}, rather than with \textit{bodies}. In the old theory, it was possible to consider a number of bodies all at the same instant, and since the time was the same for all of them it could be ignored. But now we cannot do that if we are to obtain an objective account of physical occurrences. We must mention the date at which the body is to be considered, and thus we arrive at an `event,' that is to say, something which happens at a given time. When we know the time and place of an event in one observer's system of reckoning, we can calculate its time and place according to another observer. But we must know the time as well as the place, because we can no longer ask what is its place for the new observer at the `same' time as for the old observer. There is no such thing as the `same' time for different observers, unless they are at rest relatively to each other. We need four measurements to fix a position, and four measurements fix the position of an event in space-time, not merely of a body in space. Three measurements are not enough to fix any position. That is the essence of what is meant by the substitution of space-time for space and time.
\end{quotation}
But then, unless one is prepared to accept very abstract notions which muddle the formal description of Physical Reality, in order to reconcile relativity theory with common sense, the formal requirement of a block spacetime of \textit{a priori} given events must be taken to be at the heart of Russell's recapitulation, even though he appears there only to be beating about the bush.

In fact, it is my opinion that the reason why so many such descriptions of relativistic spacetime have appeared muddled and hard to grasp, by anyone seeing the theory for the first time, is that the describers---whose point-of-departure is always to interpret the relativity of simultaneity as meaning, that two events which appear to any observer to have occurred simultaneously \textit{really} did take place simultaneously in that observer's frame; and equivalently so, for all observers, according to the principle of relativity---do not begin, before moving on to describe the way of `happenings' in spacetime in more detail, by stating that spacetime, according to the common interpretation of the mathematical formalism, is a strictly determined four-dimensional block, with the particular (somewhat peculiar, even) property that we commonly refer to as the `arrow of time'.

Now, it is not my intent to show here that this is so, because that has already been done sufficiently well by others---e.g., by those referenced above. In fact, as I \textit{will} eventually point out, in \S~\ref{sec_OIS}, it was in the basic design of the formalism originally developed by Hermann Minkowski in 1908 \cite{Einstein1952}. What I do intend to show, in Chapter~\ref{chap_DIRT}, after the more significant elements of this basic problem have been analysed in the current chapter, is that there is a \textit{more} objective interpretation of the mathematics, which is completely at odds with the original interpretation that was made, in effect, by Einstein---which he thought, incorrectly, to be the most objective one possible; for this other possibility offers a clear account of the events of spacetime in an explicitly dynamical way, which can be described through an implicit property of spacetime which is empirically supported.

The dynamical interpretation of relativity theory, which is really only a small advance on the basic idea that is described in \cite{Geroch1970}, will not be a difficult one to accept when it is shown to be entirely consistent with the mathematical theory of relativity, and is the \textit{only} interpretation of the mathematics which admits an \textit{essentially} dynamical description of the unfolding of events, because it is just the interpretation that is usually thought of anyway, even though true dynamism has not been reconciled with the theory in all essential details. Even so, it will only become clear through the basic interpretation that is described in Chapter~\ref{chap_DIRT}, precisely how it is possible that, as recently stated by George Ellis \cite{Ellis2007}, the `unchanging' `Block Universe idea, representing spacetime as a fixed whole, [which] suggests the flow of time is an illusion[, according to which] the entire universe just is, with no special meaning attached to the present time', `is best replaced by an evolving block universe [view of spacetime] which extends as time evolves, with the potential of the future continually becoming the certainty of the past; [that] spacetime itself evolves, as do the entities within it',---by explaining that even this picture should not be an accurate representation of the \textit{reality} it is used to describe, according to the most objective interpretation of the theory.

However, because the difference is purely ontological, since the mathematical description of physical events upon which it is based remains exactly the same, the argument for real dynamism shall inevitably come down to an appeal to such natural epistemological requirements as Gottfried Leibniz' `principle of sufficient reason'---viz., that there has to be a reason for everything which is sufficient to explain why it should be so, and not otherwise \cite{Leibniz1956}.\footnote{In the same paragraph in his second correspondence to Samuel Clarke, Leibniz offered two separate statements of the `principle of sufficient reason': `that nothing happens without a reason why it should be so, rather than otherwise', and `that there ought to be a sufficient reason why things should be so, and not otherwise' \cite{Leibniz1956}; the above statement is supposed to incorporate the meaning of both of these.}

And so, in order to provide an adequate argument in favour of the dynamical theory, it is imperative to investigate the essential points that have been made in the historical development of spacetime theory, because, as the above discussion has been meant to indicate, the original interpretation of the theory has had a long history before Minkowski and Einstein, and so has the one I am going to explain here, and there are many subtle points that must be highlighted before the basic justification and reasoning behind either can be completely (\textit{ergo} sufficiently) understood. For only by truly understanding the logical arguments pertaining to either theory, which are best assimilated in their nascent form, where the initial justifications can be properly dissected in order to objectively address their validity, can one free one's mind from the biased thinking that was noted by Einstein as he expressed his opinion in regard to the significance of studying not only the methodology, but also the history and philosophy of science \cite{Howard2010}:
\begin{quote}
So many people today---and even professional scientists---seem to me like somebody who has seen thousands of trees but has never seen a forest. A knowledge of the historic and philosophical background gives that kind of independence from prejudices of his generation from which most scientists are suffering. This independence created by philosophical insight is---in my opinion---the mark of distinction between a mere artisan or specialist and a real seeker after truth. 
\end{quote}

As I've said, the theory of physical Reality that I will eventually come to, incorporates the same idea of spacetime that was developed, e.g., in \cite{Geroch1970}, and basically adds to that description an unobservable noumenon; which must therefore be shown, by means that do not directly concern the observation of phenomena, to have clear physical and empirical motivation. Specifically, I will show that it is totally at odds with the very crux of the original formalism of relativistic spacetime theory, in a manner which closely resembles the advance on the Eleatic theory that was made by the Ancient Atomists---and therefore, according to the facts that may be logically deduced, is necessarily required in order to formally reconcile relativity theory with real dynamism. 

So, by making use of pertinent insights that many of the most brilliant natural philosopers in history have made, as they previously explored the Nature of Reality, while further clarifying their arguments, or advancing new arguments for or against those, by utilising the advantage of modern intelligence, a clear explanation of the present theory, which probes its basic principles, shall be presented. Then, it will be simple to see that although the theoretical description of physical observables remains exactly the same, many of the formal consequences should differ significantly from the standard conclusions that have been drawn from the original interpretation, when it has been misinterpreted as a dynamical one. As such, many paradoxes that have been previously inferred will be immediately resolved.

This concludes the introductory statement of the goals and methods of this chapter and the next. We now continue our argument by examining the theoretical achievements of Heraclitus of Ephesus, \textit{c}.~500 \textsc{b.c.}, beginning with his contribution to epistemology, which is a necessary primer to truly realising his significant contribution to natural philosophy, and is also an essential part of the method that will be used below, as it has been similarly used by many others thoughout the past two and a half millennia.
\section{On Heraclitus\label{sec_OH}}
\subsection{His Philosophy of Knowledge\label{subsec_HPN}}
As with everything else we know of Heraclitus' philosophy, his epistemology is a topic of some contention. For example, in Sextus Empiricus' \textit{Against the Logicians} \cite{Sextus2005} (126-134), the brief section on Heraclitus seems to me to begin by stating one thing about Heraclitus' epistemology, but end in showing that it is otherwise. Specifically, he begins his account with \cite{Robinson1987}, `(126) And Heraclitus -- since he \ldots supposed that man is furnished with two organs for gaining knowledge of truth, namely sensation and reason -- held \ldots that of these organs sensation is untrustworthy, and posited reason as the standard of judgment.' From everything that is now known of Heraclitus' philosophy, including the fragments that were given to us by Sextus in this very account, it seems more likely that Heraclitus did not even hold sensation and reason to be on the same level, as paths towards knowledge and understanding, but thought of sensation as the primary form, as our most direct connection to a real underlying Truth, which cannot, however, be correctly understood without clear subsequent rationalisation. As far as I can tell, this is roughly what is described by Sextus below this statement, and it is surely the message contained in the handful of other fragments that have been found elsewhere. 

It is therefore a wonder that his next statement is such a clear misconstruction of the meaning of Heraclitus' words, as, according to his own eventual conclusion, it must be regarded as grossly misleading and inconsistent. What he says, viz., is that \cite{Robinson1987}
\begin{quotation}
\noindent$<$The claim of$>$ sensation he expressly refutes with the words, `Poor witnesses for people are eyes and ears, if they possess uncomprehending (literally, `barbarian,') souls,' which is equivalent to saying, `To trust in the \textit{non-rational senses} is a mark of uncomprehending souls.'

(127) Reason, on the other hand, he declares to be the judge of truth -- not, however \textit{any} sort of reason you might care to mention but that reason which is `common' and divine. \ldots
\end{quotation}
According to what I've already said, however, and as we'll now see more clearly, the meaning of the former statement is not at all equivalent to that of the latter one. Rather, what Sextus should have said was that Heraclitus' statement is equivalent to saying, `To trust in non-rational \textit{conceptions}---those being directly subsequent to sense-perception, without having been subjected to reason---is a mark of barbarian souls.' Logically, Sextus' statement would have been more closely related to what Heraclitus wrote, if the original statement had been instead, `Poor witnesses \ldots, \textit{for} they possess \ldots'; but this is clearly not at all what Heraclitus thought, as his view was that we \textit{can} reason through to the causes of common sense, in order to discern Truth. The remainder of Sextus' argument would then have been far clearer if the sentence beginning paragraph (127) had been instead, `Reason, he declares to be the judge of truth---not, however, \textit{any} sort of reason you might care to mention, but reason \textit{about} that which is `common' and divine,'---which has a very different meaning from that which he did say. 

If it is correct to say that this is basically what Heraclitus' theory had been, then we would say that he believed in some common and divine Truth, with which we must interact in some way, in order that we might grasp it in some form, and that we should subsequently subject that which we do grasp to rational inquiry, in order to sort out what it actually, basically, is. If so, it is that which he describes as being commonly grasped, or conceived by the subconscious, that is called common sense, which must indeed occur prior to rational inspection.

In fact, this is the message that must be understood from the remainder of Sextus' account. He goes on to describe what we should now consider to be a mystical scheme of Heraclitus', for how we interact with the `common' and divine, by breathing it in, whence this comprehending reason is given its divine birth, and eventually concludes \cite{Robinson1987},
\begin{quotation}
(131) Heraclitus says, then, that this `common' and divine \textit{logos} -- by participation in which we become rational -- is the yardstick of truth. So that which appears $<$such and such$>$ to all in common is trustworthy (for it is grasped by the `common' and divine \textit{logos}), but that which strikes only an individual as $<$such-and-such$>$ is -- for the opposite reason -- \textit{un}trustworthy.

(132) Thus the aforementioned man begins his work \textit{On Nature}, and in a certain fashion points out $<$the existence of$>$ the surrounding $<$nature$>$ with the words:\footnote{I've chosen to use Jonathan Barnes' translation of this fragment, from \cite{Barnes1982}, rather than the one from \cite{Robinson1987}, because it retains the ambiguity of the application of `always' here, which is commonly agreed upon to have been deliberately constructed, so that a comma could be thought to be placed before or after `always' in order to mean either `which is the case always', or `always men prove to be uncomprehending [of it]'; whereas the latter translator, who agrees with this interpretation, has nevertheless explicitly written `forever' twice. Furthermore, it seems that `always' accounts better for the original meaning of the statement, by expressing a generality of ways, than does `forever', which is strictly temporal.} `And of this account (\textit{logos}) which is the case always men prove to be uncomprehending, both before they hear it and once they have heard it. For although everything comes about in accordance with this account (\textit{logos}), they are like inexperienced men when they experience both the words and the deeds of the sort which I recount by dividing up each thing in accordance with its nature (\textit{phusis}) and saying how it is; but other men do not notice what they do when they are awake, just as they are oblivious of things when asleep.'

(133) For having in these words expressly stated the view that we do and apprehend everything thanks to our participation in the divine \textit{logos}, he goes on a little further, then adds: `That is why it is necessary to follow that which is $<$common$>$. Though the account (\textit{logos}) is common, however, the many live as though they had a private understanding'. This \textit{logos} is nothing else than an explanation (exposition, articulation, \textit{ex\={e}g\={e}sis}) of the mode of arrangement of the universe. That is why we speak truly whenever and in so far as we share in the recollection of it (ie, of the \textit{logos}) but are invariably mistaken on matters of private opinion.

(134) So here and in these words he states clearly that the common \textit{logos} is the yardstick $<$of truth$>$; the things that appear such and such \textit{in common} are trustworthy, as being judged by the common \textit{logos}, whereas those that appear $<$such and such$>$ to each person privately are false.
\end{quotation}
So he ends his account of Heraclitus' doctrine, basically by saying that he had believed in the primacy of common sense, which we must trust to, in order to reason correctly, and only by rationalisation that is faithful to that common \textit{logos}, can we come to a better understanding of Nature's Truth---i.e., expose the True mode of arrangement of the Universe, which he has already assumed to have a prior Real existence.

In summary, this account is understood to say that according to Heraclitus, common sense must be used in order to realise what is True, but its origin \textit{in} Truth cannot be understood unless it is considered with divine reason; i.e., rational, insightful reason, that is always held to be consistent with what is common, as it probes its own very roots. In other words, observations and logical argument must be used in order to find the prior causes of that which everyone knows through common sense. In essence, Heraclitus placed common sense prior to logical inquiry, for it must be the result of something implicit in Nature, but he says that only those who think carefully about it will come to any clear understanding, whereas barbaric misunderstandings result from being careless. In this way, Heraclitus began his account \textit{On Nature} with a direct challenge to all who seek to understand It, that they must always follow this path, for it is the only one through which Truth can be properly exposed.

Many of the other ($\sim130$) remaining Heraclitean fragments attest to this reconstruction of Sextus' account \cite{Robinson1987}:\footnote{The numbers given here are standards.}
\begin{quote}
17. Many people do \textit{not} `understand the sorts of thing they encounter'! Nor do they recognise them $<$even$>$ after they have had experience $<$of them$>$ -- though they themselves think $<$they recognize them$>$.

34. Uncomprehending $<$even$>$ when they have heard $<$the truth about things?$>$, they are like deaf people. The saying `absent while present' fits them well (literally, `bears witness to them').

47. Let us not make random conjectures about the most important matters.

95, 109. It is better, [says Heraclitus,] to conceal ignorance.

112. Sound thinking $<$is$>$ a very great virtue, and $<$practical$>$ wisdom $<$consists in our$>$ saying what is true and acting in accordance with $<$the$>$ real constitution $<$of things$>$, $<$by$>$ paying heed $<$to it$>$.

113. Thinking is common to all. [This is given in conjunction with,]

114. Those who $<$would$>$ speak with insight must base themselves firmly on that which is common to all, as city does upon $<$its$>$ law -- and much \textit{more} firmly! For all human laws are nourished by one $<$law$>$, the divine $<$law$>$. For it holds sway to the extent that it wishes, and suffices for all, and is still left over.

116. All people have a claim to self-knowledge (literally, `self-ascertainment') and sound thinking.
\end{quote}
And more specific to sense-perception,
\begin{quote}
28a. The most esteemed $<$of people$>$ `ascertains' -- and holds fast to! -- what $<$merely$>$ seems $<$to be the case$>$.

46. [He used to say that] thinking is $<$an instance of the$>$ sacred disease [and that] sight is deceptive.

54. An unapparent connection is stronger (or: better) than one which is obvious.

55. Whatsoever things $<$are$>$ objects of sight, hearing, $<$and$>$ experience -- these things I hold in higher esteem.

72. They are separated from that with which they are in the most continuous contact.

101a. Eyes are more accurate witnesses than are [the] ears.

123. $<$a thing's? (the world's?)$>$ real constitution [according to Heraclitus] has a tendency to conceal itself.
\end{quote}

This is what we will take Heraclitus' epistemology to have been, which is in total agreement with the stance I have taken in the present analysis, as I have attempted to state in the introduction. Accordingly, we will have no more use for any theory reached by the pure reasoning based on any set of facts, which is not also basically consistent with common sense, than we will have for a common sense-based conclusion that is arrived at without sufficient reason. As such, our main problem will be, as it always has, and may always be,---that we must \textit{only} trust in the completeness of our rationalisations \textit{until} they should show themselves as having been objectively inconsistent, or lacking sufficient reason. However, we find that the surest way to avoid error, as Heraclitus said in so many words, is to use rational insight in order to derive clear expressions of the principles from which our theories are built; for only by clearly recognising what those basic principles are, which must be at once consistent with common sense and all other observations, can we truly understand our theories enough to be capable of objectively assessing their strengths and weeknesses. In other words, his epistemology is based upon a single principle, that all theories come from notions that originate in observation, which bears witness, in some fashion, to an underlying Truth; but in order for any theory to be considered valid (though ultimately refutable, as we can never observe, nor think of, all possibilities), and to be essentially verifiable, those notions must first become rationalised principles.

Historically, Heraclitus has been far from unique in his epistemological stance, which was later adopted by many Ancient schools of thought. For example, we see this in Galen's report of Democritus' theory \cite{Taylor1999}:
\begin{quotation}
If someone cannot even make a start except from something evident, how can he be relied on when he attacks his very starting-point? Democritus was aware of this; when he was attacking the appearances with the words `By convention colour, by convention sweet, by convention bitter, but in reality atoms and void' he made the senses reply to thought as follows: `Wretched mind, you get your evidence from us, and yet you overthrow us? The overthrow is a fall for you'.
\end{quotation}
to which the same commentator elaborates further elsewhere \cite{Taylor1999}:
\begin{quote}
By the expression `by convention' he means `conventionally' and `relative to us,' not according to the nature of things themselves, which he calls by contrast `reality,' forming the term from `real' which means `true.'
\end{quote}
According to Diogenes La\"{e}rtius, Democritus is also to have said that `In reality we know nothing, for truth is in the depths' \cite{Taylor1999}; but we know that Democritus did apply reason to the sensible data, in his life's attempt to discover what he could, of that deep truth, with greater tenacity than any who had come before him; for although he realised how elusive it was, all accounts tell that he spent all of his inheritance travelling the world in search of knowledge, and, even, according to Eusebius, that he was reported to have said that he `would rather discover a single explanation than acquire the kingdom of the Persians' \cite{Taylor1999}; and Aristotle---though his beliefs and his theory of Nature had been quite different---praised Democritus' use of this method, in \textit{On Generation and Corruption} \cite{Aristotle1984}: `In general, no one except Democritus has applied himself to any of these matters [sc., coming-to-be and passing-away] in a more than superficial way. Democritus, however, does seem not only to have thought about all the problems, but also to be distinguished from the outset by his method.' A little further down, he writes that 
\begin{quotation}
\ldots those who dwell in intimate association with nature and its phenomena are more able to lay down principles such as to admit of a wide and coherent development; while those whom devotion to abstract discussions has rendered unobservant of the facts are too ready to dogmatize on the basis of a few observations. The rival treatments of the subject now before us will serve to illustrate how great is the difference between a scientific and a dialectical method of inquiry. For, whereas the one school argues that there must be atomic magnitudes because otherwise The Triangle will be more than one, Democritus would appear to have been convinced by arguments appropriate to the subject, i.e.~drawn from the science of nature;
\end{quotation} 
and, eventually, when he describes the atomic theory, he begins by saying that `The most systematic theory \ldots, and one that applied to all bodies, was advanced by Leucippus and Democritus: and, in maintaining it, they took as their starting-point what naturally comes first.' So we see that Democritus, and also Aristotle, as we'll see in greater detail in \S~\ref{sec_ANE}, saw great worth in the practice of this Heraclitean form of natural epistemology.

It should also be admitted, however, that Sextus proceeds directly from his bastardisation of Heraclitus' philosophy, in \cite{Sextus2005}, to similarly misconstrue the meaning in Democritus' writing, according to a few selected quotations; e.g., he quotes a similar passage as Galen has, that `By convention \ldots', along with similar other ones, but leaves out the senses' reply. Thus, it appears he either completely mistook, or deliberately misconstrued the meaning that had been expressed by both of these men: that primary conceptions are untrustworthy accounts of True perceptions, which can only be correctly understood through subsequent reasoning, but are nevertheless patently essential to the process of true knowledge acquisition. In other words, although, as Sextus tells us, both Heraclitus and Democritus had written that sense-perception and reason are our two inherent means of ascertaining Natural Truths (which they both believed in), and both considered sense-perception to be untrustworthy, whereas reason they thought to be reliable, the point that both had obviously tried to make, really comes down to the fact that perceptions themselves are our basic data of observation, and that biased interpretations made by the subconscious are untrustworthy, so that subsequent reasoning is necessary to form a better understanding of what those data Really are. 

Probably the clearest statement of the use of this method in the development of any Ancient theory, however, which no Sceptic could deny, is given in the surviving writings of the Epicureans, who resurrected the atomic theory, in Athens, roughly fifteen years after Aristotle's death. Although most of these works were lost, a few important summaries were preserved by Diogenes La\"{e}rtius. In particular, Epicurus' letter to Herodotus \cite{Diogenes1931} (X.~35-83)---his surviving epistle on physics---which was prepared as an `epitome and manual of the doctrines as a whole', now exists, through Diogenes, in its entirety. There, Epicurus states that `we must by all means stick to our sensations, that is, simply to the present impressions whether of the mind or of any criterion whatever, and similarly to our actual feelings, in order that we may have the means of determining that which needs confirmation and that which is obscure'; for, `it is upon sensation that reason must rely when it attempts to infer the unknown from the known.'

As well, Diogenes preserved a list of forty `Sovran Maxims', composed by the Epicureans from the writings of Epicurus, of which three are relevant here \cite{Diogenes1931}:
\begin{itemize}
\item[22.]{`We must take into account as the end all that really exists and all clear evidence of sense to which we refer our opinions; for otherwise everything will be full of uncertainty and confusion.'}
\item[23.]{`If you fight against all your sensations, you will have no standard to which to refer to, and thus no means of judging even those judgements which you pronounce false.'}
\item[24.]{`If you reject absolutely any single sensation without stopping to discriminate with respect to that which awaits confirmation between matter of opinion and that which is already present, whether in sensation or in feelings or in any presentative perception of the mind, you will throw into confusion even the rest of your sensations by your groundless belief and so you will be rejecting the standard of truth altogether. If in your ideas based upon opinion you hastily affirm as true all that awaits confirmation as well as that which does not, you will not escape error, as you will be maintaining complete ambiguity whenever it is a case of judging between right and wrong opinion.'}
\end{itemize}
Or, according a recent gloss of the Epicurean epistemology by Robert Hicks \cite{Hicks1962}: `It is only through sense that we come into contact with reality; hence all our sensations are witnesses to reality. The senses cannot be deceived. There can be no such thing, properly speaking, as sense-illusion or hallucination. The mistake lies in the misinterpretation of our sensations. What we suppose that we perceive is too often our own mental presupposition, our own over-hasty inference from what we actually do perceive. \ldots. Sensations themselves must be scrutinised, and the element which the mind itself has added must be removed before we get back to the original data, the perceptions which put us in touch with reality.'

Similarly, the Stoic epistemology, which was related by Diogenes in his account of their philosophy, is thought to have risen from the foundation laid by Heraclitus. Basically, their goal was to form a consistent understanding of the grand spectrum of things by understanding `presentations', or mental impressions, which are imprints on the soul, through subsequent thought \cite{Diogenes1931}. As Diogenes related for us from an earlier text \cite{Diogenes1931} (VII.~49),
\begin{quotation}
``The Stoics agree to put in the forefront the doctrine of presentation and sensation, inasmuch as the standard by which the truth of things is tested is generically a presentation, and again the theory of assent and that of apprehension and thought, which precedes all the rest, cannot be stated apart from presentation. For presentation comes first; then thought, which is capable of expressing itself, puts into the form of a proposition that which the subject receives from a presentation.''
\end{quotation} 
Stoicism was founded by Zeno of Citium, also in Athens, shortly after Epicurus founded his school there, and its doctrine was widely followed for more than eight hundred years. Regarding its founder, Hicks writes \cite{Hicks1962}, `To Zeno \ldots, the natural was the rational and the first mark of reason was self-consistency. A rational life must follow a single harmonious plan, whereas the paths of folly are many and various, but always stamped with inconsistency and contradiction. In this postulate of a rational law the Stoics had a precursor in Heraclitus.'
\subsection{His Natural Philosophy\label{subsec_HNP}}
So much for Heraclitus' philosophy of knowledge and its ancient legacy: let us now turn to his philosophy of Nature---or, rather, a specific aspect of it which is the most important result of his having directly applied this method in order to try to understand Nature, which was indeed an essential step towards the present theory of Time. To be clear, most of what is known about Heraclitus' natural philosophy has to do with his mentioning of the unity of contraries, and he held that the primary element of Nature was fire. It is not immediately relevant to delve into any of this but to say that it is reasonable to admit, as, in my opinion, Barnes' argument \cite{Barnes1982} sufficiently shows, that Heraclitus' thesis was based on the principle that Nature, as a whole, is uniformly and continuously in flux---viz., in Time.---And the first explicit recognition of this principle, among the natural philosophers, lies with Heraclitus; for, if we are to take Barnes' word \cite{Barnes1982}, even
\begin{quote}
Some of those scholars who accept the Theory as Heraclitean are inclined to see nothing very original in it: the Milesians, after all, had held a similar view. The Milesians, like all observant men before Parmenides, had indeed noticed that things change: the world is patently not a static \textit{tableau}. Yet it is far from a patent truth that \textit{everything} changes, still less that everything \textit{always} changes; and the Milesians, like ordinary men before Heraclitus, seem to have thought that within the changing world there was room for a number of stable courses, and the earth does not move from its place. There is no reason to deny Heraclitus the novelty of generalizing the natural view of a changing world to the more pugnacious thesis that everything changes\ldots
\end{quote}

Indeed, the common notion that everything progresses through time cannot---by definition!---be attributed to Heraclitus; e.g., we read in Plato's \textit{Cratylus} (402a-b) \cite{Plato1997},
\begin{quotation}
\textsc{Socrates}: I seem to see Heraclitus spouting some ancient bits of wisdom that Homer also tells us---wisdom as old as the days of Cronus and Rhea.

\textsc{Hermogenes}: What are you referring to?

\textsc{Socrates}: Heraclitus says somewhere that ``everything gives way and nothing stands fast,'' and, likening the things that are to the flowing ($rho\overline{e}$) of a river, he says that ``you cannot step into the same river twice.''

\textsc{Hermogenes}: So he does.

\textsc{Socrates}: Well, then, don't you think that whoever gave the names `Rhea' and `Cronus' to the ancestors of the other gods understood things in the same way as Heraclitus? Or do you think he gave them both the names of streams (\textit{rheumata}) merely by chance?\symbolfootnote[1]{`Rhea' sounds a lot like `\textit{rheuma}' (`stream'); apparently Socrates expects Hermogenes to hear `Cronus' as connected with `\textit{krounos}' (`spring'). [Translator's note.]} Similarly, Homer speaks of `Ocean, origin of the gods, and their mother Tethys'; I think Hesiod says much the same. Orpheus, too, says somewhere that `Fair-flowing Ocean was the first to marry, and he wedded his sister, the daughter of his mother.' See how they agree with each other, and how they all lean towards the doctrines of Heraclitus.
\end{quotation}

But if we understand Heraclitus' epistemology correctly, we must recognise the significance of his own achievement of rational insight, to go beyond that which everyone before him knew, and clearly realise that \textit{a priori} everything continually progresses uniformly through Time. Although the subconscious notion, we say, must come from the very essence of Nature, and has therefore always been, and must be, common among sentient beings, it is truly fair among humans---do with it what we will---to credit the first clear recognition of this fact to Heraclitus.

However, we must do full justice to the true brilliance of Heraclitus' insight, because the common belief among modern philosophers is that Plato did not. Namely, it is now understood that the river fragments, in which Heraclitus' principle is considered most apparent, are not supposed to be thought of as analogical or metaphorical---i.e., Heraclitus is not likening the change that occurs, as Time continuously progresses, to the flow of a river, as Plato says; nor is he even using the flow of the river symbolically, as a representation of the flow of the Universe through Time---but only as an example which clearly illustrates the uniform flux of all matter that exists throughout the Universe, regardless of its state of motion.

Indeed, consider three fragments \cite{Robinson1987}:
\begin{quote}
84a. While changing it rests.

30. $<$The ordered?$>$ world, the same for all, no god or man made, but it always was, is, and will be, an everliving fire, being kindled in measures and being put out in measures.

12. As they step into the same rivers, different and $<$still$>$ different waters flow upon them.
\end{quote}
Only a further brief note to accompany each of these will be necessary:---in fragment 84a, the awkward ordering of the two verbs impresses the primacy of change, even in the thing that rests, as rest requires existence as much as any other state of being, and therefore an interval of change;---some critics have argued that the literal meanings of the surviving Heraclitean fragments can all be construed so as to admit no evidence of a principle of flux---but any critic who has ever considered that any part of Nature, according to Heraclitus' theory, might not be in continuous flux, can surely never have seen a fire---a singular thing that is never the same anywhere, from one moment to the next;---there are actually three river fragments, but the one given here is considered to be the primary one; it must be interpreted as a singlular statement regarding two different things: viz., that in one sense, humans and rivers always remain the same things; but, in another sense, that both of them continuously and uniformly change, so that although they always exist together, both do change from one instant to the next; and further, this is implicit in the fact that they \textit{remain} the same things, which is yet another state of being.

Perhaps one of the clearest explanations of Heraclitus' principle, which nevertheless begins with the common connection to flowing water, using it metaphorically, although eventually placing another of the river fragments in context appropriately, comes from Seneca the Younger's \textit{Ad Lucilium Epistulae Morales} (LXIII, 23) \cite{Seneca1917}:
\begin{quotation}
Our bodies are hurried along like flowing waters; every visible object accompanies time in its flight; of the things which we see, nothing is fixed. Even I \ldots, as I comment on this change, am changed myself. This is just what Heraclitus says: ``We go down twice into the same river, and yet [each time] into a different river.''\footnote{`In idem flumen bis descendimus et non descendimus';---more literally, `We twice descend and do not descend into the same river.' The original Latin is therefore more ambiguous than this translation suggests, as it \textit{may} be taken to mean that the two existences apply singularly to the act of descent, rather than to the different states of the physical system. However, given the context, we can be sure that \textit{Seneca's} meaning has been translated correctly, as he probably did not mean to state something that has nothing whatsoever to do with his line of thought. I have therefore added the words `each time', in order to give further clarification to this interpretation.} For the stream still keeps the name, but the water has already flowed past. Of course this is much more evident in rivers than in human beings. Still, we mortals are also carried past in no less speedy a course; \ldots every instant means the death of our previous condition. \ldots So much for a man,---a substance that follows away and falls, exposed to every influence; but the universe, too, immortal and enduring as it is, changes and never remains the same. For though it has within itself all that it has had, it has it in a different way from that in which it has had it; it keeps changing its arrangement. 
\end{quotation}
Thus, according to Seneca's interpretation of Heraclitus' statement, it is not the flow of the river that is to be thought of as being similar to time, but that the idea of a flowing river serves to illustrate more clearly that which is the same thing implicit in humans and rivers, as it is in all of Nature, but which is less obvious in an inert body,---that Time continuously proceeds, so that every instant is essentially different from all others---every instant in the continuous stream of Time means the death of all the Universe's previous condition, and the birth of a new one, in which the Universe nevertheless remains the unity of all things material.

Heraclitus was therefore the first true relativist:---who first realised that no rest-frame is Truly static;---who realised, before all later cosmologists, that there is a Universal rest-frame---a mean to which all the heavenly bodies and the flowing rivers within them, exist in relative motion---yet even the bodies that remain at rest with respect to that frame continually change, as the Universe changes---though while changing, It rests. 
\section{On Eleatism\label{sec_OE}}
As we now move from Heraclitus', to a description of the Eleatic theory, it is worth noting, first of all, that, whereas Heraclitus' theory was based on the reduction of observational evidence to basic principles, the Eleatics based theirs on systemmatic logical deduction. Now, our primary reason for analysing the Eleatic theory, is so that we may investigate an extremely serious philosophical problem that was first propounded by Parmenides, which remains a significant problem in understanding the meaning of relativity theory; but because it has general relevance in our present endeavours, and so that we may clearly understand their theory---in particular, the origin of this persisting problem,---we also have cause to further investigate their method of inquiry. 

The best way to understand the fragmentary statement of Parmenides' theory is to begin with a clear description of the basic principle of his logical deduction. The reason for this, is that in its fragmentary form Parmenides' poem is probably too ambiguous for a conclusive literary deduction, but it seems that with this prior understanding, its meaning immediately becomes quite clear, as the reasoning behind each line can be more easily grasped.

The general idea is that what appears to be a certain conclusion that could be taken from Parmenides' argument, is actually the prime motivation for his entire theory of nature, which results from a straightforward deduction based on a single principle that he decided to assume, in all likelihood incorrectly, and which is corroborated by the resulting theory's partial resolution of the problem of becoming. More specifically, consider a younger Parmenides, thinking about Nature---about being, becoming, and change:---he reasons that Being Is and Not-being Isn't, and nothing which Isn't can actually exist. If Change Is,---from what, and where to? If from, and into, That-which-does-not-presently-exist, then surely Change cannot Be, because past-state could not thereby become What-is, nor can present-state become What-is-not. Therefore, the past and future, which common sense tells us were, and will be, must Be \textit{now}, as with the present, and time must be an illusion. Therefore, the One that is All must Be all of space and time, absolutely determined in all four dimensions, in a singular temporal sense, as a Block.

The remainder of the theory, which is brilliant in its own right, is built upon this; I'll remark only on one aspect which should be of interest to the modern relativist: viz.~that Parmenides realised,---as Einstein also would, two and a half millennia later, for his own particular reasons,---that there would be a problem with his cosmology if space would be infinite; so he did \textit{precisely} what Einstein did, and made it spherical. However, so that it would have no beginning or end in time, he made the One (what can \textit{now} only be recognised as) a 4-sphere (although his theory was not formally so explicit), which he explained by the only conceptually possible analogy (in any millennium) of an embedded 2-sphere.

Now, what should become plain through consideration of the above interpretation of what Parmenides' logic \textit{seems} to have been, is that his entire theory hinges on one basic principle: that `time' and `space' have no real existence in and of themselves, independent of the material that Is there; i.e., that the One is \textit{purely} material, and there is no `where' for anything to Be---that Being is \textit{material} and Not-being is \textit{immaterial}, and \textit{matter} is all that there Is. As Giorgio de Santillana says \cite{deSantillana1961}, the concept of Being in Parmenides' poem is pure geometrical space, and `[g]eometry as the Greeks meant it put three requirements on its space: first, it must have continuity (in a sense somewhat stronger than the mere absense of gaps between points); second, it must be the same, homogeneous throughout, so that we can move figures freely from place to place without altering their geometrical properties; and, finally, it must be isotropic, or the same in all directions'; furthermore, it is a non-trivial fact that he does actually formulate this as a purely \textit{corporeal} existence, with matter and existence being perfectly synonymous things, and all that there Is, so that the concept of `place', e.g., as defined later by Aristotle, in his \textit{Physics} \cite{Aristotle1984}, is purely a delusion; i.e., that the Reality that Is is to be a purely physical Corporeality, existing as an isotropic, homogeneous, complete metric space, though somehow also abstract to the senses. 

If this basic principle had actually been correct, it would be much more difficult to disagree with Parmenides' theory than we shall presently find, as Plato and Aristotle both realised, who both believed that nothing physical could be incorporeal, which problem they circumvented in proposing the reality of abstract, nonphysical things. This problem, first raised by the Eleatics, is of paramount importance in metaphysics; therefore, we must take a brief pause, in order to examine the various forms of real existence that were proposed as a direct result of the Eleatic reasoning, so that we may better understand both the choice they made and, eventually, our own metaphysical theory.

We therefore begin by distinguishing the classifications of corporeal and incorporeal, and physical and nonphysical, as defining absolute complements of two different forms of real existence. In modern physics, e.g., we say that the elementary particles and the compounds formed from them are the corporeal things, whereas the vacuum of spacetime is incorporeal; and we say that anything---i.e., any body, process, state, etc.---that is subject to the fundamental Laws of physics, whatever they may be, is physical, whereas anything else is nonphysical. In fact, it remains a possibility that all nonphysical things might have physical counterparts; e.g., that an idea should accompany a certain physical situation and occurrence of synapses in a brain. However, it is not our purpose here to judge whether any nonphysical thing could actually be real, in and of itself, or whether all must be byproducts as such; rather, in our later theoretical development, we shall stick to the forms of physical existence, with the aim of giving a sufficient self-consistent description of them. 

But before we can assume that, we must see how the various theories of real existence generally apply in theories of Nature. First of all, it may be reasonably posited that all things corporeal must be physical; so we are led to identify three possible forms of real existence: physical corporeal, physical incorporeal, and nonphysical incorporeal. Then, as I've said in so many words, the Eleatics believed, and this was the principal tenet of their theory, that the only real existence could be a physical corporeal continuous geometry; to them, whether a body's \textit{place} was a physical incorporeal container, or a useful concept for describing things in common speach, as they exist in the mind, did not matter because neither would actually be real. And if `space' and `time' are not real things, then what Is could not actually travel through them---it could not come from the past, or move into the future, any more than it could move from place to place. 

Following the Eleatic line of reasoning, but not their ultimate beliefs, the Atomists, Leucippus and Democritus, posited that a physical and incorporeal void should also be real, though they continued to deny the actual existence of anything nonphysical. In their theory, the void would be continuous and would contain matter composed of finite indivisible atoms. Conversely, Plato took up the alternative,---that reality consists of the physical corporeal and the nonphysical incorporeal,---and Aristotle further developed the physical theory that would naturally describe this, i.e., eventually arriving at an abstract physical description of time, in which the material Universe would be continually rearranged; for so he found that matter should be everywhere, but that the places containing matter could also be real, and change as well; and that acting by mutual contact, `some things will suffer action and others will act, provided they are by nature adapted for reciprocal action and passion' \cite{Aristotle1984}. As we'll see in \S~\ref{sec_OT}, it was St. Augustine who eventually provided the sufficient argument that favours Time as a physical thing, though he did not himself believe this.

Thus, I think---and this is precisely how the problem was taken up later by the Atomists, whose methods in this respect Aristotle placed above all others, as mentioned above---that the problem with Parmenides' thesis is merely that he started off on the wrong foot: he began from a basic principle he believed was reasonable to assume, then used further reason and reached an inconsistency; but rather than re-tracing his steps and objectively inspecting his basic principles, he forced himself to adjust his understanding of Nature, describing it as an abstract and unverifiable thing, in an attempt to reconcile the issues he had found---and he was probably happy to do so, as he found Genesis to be less of a problem. Personally, if I wished to expound such a theory with artful craft, I should hope to have written something like this \cite{Gallop1984}:
\begin{quotation}
\noindent Come, I shall tell you, and do you listen and convey the story,\\
What routes of inquiry alone there are for thinking:\\
The one---that [\textit{it}] \textit{is}, and that [\textit{it}] \textit{cannot not be},\\
Is the path of Persuasion (for it attends upon truth);\\
The other---that [\textit{it}] \textit{is not} and that [\textit{it}] \textit{needs must not be},\\
That I point out to you to be a path wholly unlearnable,\\
For you could not know what-is-not (for that is not feasable),\\
Nor could you point it out.\\
\ldots\\
It must be that what is there for speaking and thinking of \textit{is};~for [it] is $\mathrm{there\!~to\!~be,}$\\
Whereas nothing is not; that is what I bid you consider,\\
For $<$I restrain$>$ you from that first route of inquiry,\\
And then also from this one, on which mortals knowing nothing\\
Wander, two-headed; for helplessness in their\\
Breasts guides their distracted mind; and they are carried\\
Deaf and blind alike, dazed, uncritical tribes,\\
By whom being and not-being have been thought both the same\\
And not the same; and the path of all is backward-turning.\\
\ldots\\
For never shall this prevail, that things that are not \textit{are};\\
But do you restrain your thought from this route of inquiry,\\
Nor let habit force you, along this route of much-experience,\\
To ply an aimless eye and ringing ear\\
And tongue; but judge by reasoning the very contentious disproof\\
That has been uttered by me.\\
\ldots

A single story of a route still\\
Is left: that [\textit{it}] \textit{is}; on this [route] there are signs\\
Very numerous: that what-is is ungenerated and imperishable;\\
Whole, single-limbed, steadfast, and complete;\\
Nor was [it] once, nor will [it] be, since [it] is, now, all together,\\
One, continuous; for what coming-to-be of it will you seek?\\
In what way, whence, did [it] grow? Neither from what-is-not shall I allow\\
You to say or think; For it is not to be said or thought\\
That [\textit{it}] \textit{is not}. And what need could have impelled it to grow\\
Later or sooner, if it began from nothing?\\
Thus [it] must either be completely or not at all.\\
Nor will the strength of trust ever allow anything to come-to-be from what-is\\
Besides it; therefore neither [its] coming-to-be\\
Nor [its] perishing has Justice allowed, relaxing her shackles,\\
But she holds [it] fast; the decision about these matters depends on this:\\
\textit{Is} [\textit{it}] \textit{or is} [\textit{it}] \textit{not}? but it has been decided, as is necessary,\\
To let go the one as unthinkable, unnameable (for it is no true\\
Route), but to allow the other, so that it is, and is true.\\
And how could what-is be in the future; and how could [it] come-to-be?\\
For if [it] came-to-be, [it] is not, nor [is it] if at some time [it] is going to be.\\
Thus, coming-to-be is extinguished and perishing not to be heard of.\\
Nor is [it] divisible, since [it] all alike \textit{is};\\
Nor is [it] somewhat more here, which would keep it from holding together,\\
Nor is [it] somewhat less, but [it] is all full of what-is.\\
Therefore [it] is all continuous; for what-is is in contact with what-is.\\
Moreover, changeless in the limits of great chains\\ 
\hspace{0pt}[It] is un-beginning and unceasing, since coming-to-be and perishing\\
Have been driven far off, and true trust has thrust them out.\\
Remaining the same and in the same, [it] lies by itself\\
And remains thus firmly in place; for strong Necessity\\
Holds [it] fast in chains of a limit, which fences it about.\\
Wherefore it is not right for what-is to be incomplete;\\
For [it] is not lacking; but if [it] were, [it] would lack everything.\\
The same thing is for thinking and [is] that there is thought;\\
For not without what-is, on which [it] depends, having been declared,\\
Will you find thinking; for nothing else $<$either$>$ is or will be\\
Besides what-is, since it was just this that Fate did shackle\\
To be whole and changeless; wherefore it has been named all things\\
That mortals have established, trusting them to be true,\\
To come-to-be and to perish, to be and not to be,\\
And to shift place and to exchange bright colour.\\
Since, then, there is a furthest limit, [it] is completed,\\
From every direction like the bulk of a well-rounded sphere,\\
Everywhere from the centre equally matched; for [it] must not be any larger\\
Or any smaller here or there;\\
For neither is there what-is-not, which could stop it from reaching\\ 
\hspace{0pt}[Its] like; nor is there a way in which what-is could be\\
More here and less there, since [it] all inviolably \textit{is};\\
For equal to itself from every direction, [it] lies uniformly within its limits.\\
Here I stop my trustworthy speech to you and thought\\
About truth; from here onwards learn mortal beliefs,\\
Listening to the deceitful ordering of my words \ldots
\end{quotation}

As Barnes points out \cite{Barnes1982}, Parmenides' doctrine is not based on an epistemological principle that common sense must be denied in objective philosophical deduction,---that he does not say, `don't listen merely to your eyes and ears, and bear witness to others' testimonies which are based on sense-perception'; but says `listen to reason, and if reason produces a result that is objectionable to common sense, do not deny it without a reasonable rebuttal'.

This is a reasonable enough assertion; but it is equally reasonable to say that if some observation of Nature poses a problem to a theory that has been deduced from a certain set of facts, the worst kind of cop-out for the natural philosopher is to immediately assume that thing is not really a part of Nature---that Nature must be otherwise---and then attempt to construct a philosophy based on that assumption, without first objectively examining the principles which led to the original problem---for many more shall inevitably follow. Here, Parmenides seems to have disposed of two great philosophical problems by denying the common sense notions of change, alteration, flux, motion, etc.---but only provisionally, as he has done this without having adequately addressed the more critical problem of reconciling his theory with the empirical evidence. 

Here, I say `\textit{seems} to have disposed', because I wonder,---has he really done away with the problem of the origin of Nature by assuming all four dimensions are really a temporally singular entity? He has not explained the origin of that singularity, but he argues correctly that this mode of inquiry is a misunderstanding of his thesis. So we ask instead, Why \textit{is} reality the One? To which we find only the entirely unacceptable anthropic or divine responses. I concede that his may be closer to a solution than, say, spontaneous generation, which is met with the multiple snags well-known in Big Bang cosmology; but it is no better, in this regard, than what the Atomists later came up with, in response to Parmenides' argument. So we should in fact say that he has \textit{circumvented} the problems of the Universe's coming-to-be, and why It is as it Is.

There is an important contrast that can be made with Galileo's achievement in discovering the relativity of inertia: i.e., that when one principle, which is necessarily based on some empirical evidence, appears to be in conflict with some other known datum, it cannot do to simply deny either of these, rather than examining all of the data objectively in order to find a principle consistent with all, as it could well be the principle in conflict that is incorrect; i.e., that the basic principles used in the deduction of a theory must be reduced when that theory has become liable---and it is not surprising that this was indeed the problem with the first theory ever logically deduced. 

But in Galileo's case, the principle that the Earth is at the centre of the Universe was based on the fact that we immediately infer this from basic perception, and had therefore been part of many theories since Ancient times---though not all of them. However, observations of the motions of celestial bodies, and the principle of a True physical order in Nature, could not be reconciled with the geocentric principle, which similarly led Kepler to the discovery of his Three Laws---although Galileo, who probably didn't know of Kepler's Laws \cite{deSantillana1955}, for it seems that Kepler himself never correctly understood the worth of his \textit{physically meaningful} discoveries, and therefore never supported them as he did his more fantastic endeavours \cite{Koestler1959}, based his argument for heliocentrism primarily on evidence gained from telescopic observations \cite{Galileo1967}. Therefore, the problem was met by two conflicting principles, each supported by different empirical evidence, and the problem became to determine which was false. By reasoning that we would not actually notice any difference if the Earth were in a heliocentric orbit, due to the relativity of inertia, which is also empirically supported, and shows that the common sense-based geocentric principle was basically flawed, Galileo was able to advance an intuitive argument that was again supported by empirical evidence.

Conversely, Parmenides gave no empirical evidence to support his claim that we are delusional about change and motion; his argument was that this should have been accepted because it was the result of a more substantial, logically deduced theory, and that the need for general consistency should be trumped by such pure reason, even if that should neglect some relevant facts. In fact, the very resolution to his problem lies in that which he has led himself to deny: for if we move through space and time, as common sense tells us we do, then space and time should be ontologically real, regardless of whether something is actually there. The logical implication is that there \textit{could be} places in space or time where no material is present,---that spacetime is not a continuous material \textit{plenum},---which would then allow for the motion of matter through those void regions. Although the Eleatics apparently did realise this, they abhorred the prospect that Not-being could Be, and thus failed to accept the possibility of separating the incorporeal from the corporeal within the realm of physical existence---and this is precisely what Leucippus and Democritus \textit{did} do.

However, it is obvious, and has always been clearly understood, that the Eleatic doctrine cannot be refuted scientifically; for nothing can be ultimately refuted if it cannot possibly be detected, and this is indeed the case with the Eleatic hypothesis that there is one singular eternity. Science inductively infers hypotheses and verifies their correctness through as many different repeatable observations as possible, or else falsifies them by a single one; therefore, an hypothesis that is impossible to verify \textit{once} is not scientific. But one which is totally at odds with the data we do possess, however untrustworthy we might think they are, is also physically useless. The truth is, that an intuitive example that would suggest Heraclitus' principle of Universal flux---the first principle of relativity---is incorrect, has never been found in Nature---and therein lies the best scientific evidence we can ever hope to find, that it is. 

Generally speaking, without a sufficient reductive argument to provide a new, but nevertheless equally acceptable, intuitive understanding of Nature, any theory that may be logically deduced, which finds itself in stark conflict with the appearance of Reality, and the principles that may be directly inferred from that, cannot be allowed to replace such natural intuition; its merits may be weighed, and, even if it is found to have predictive power beyond that of any other contending theory, it shall never be an acceptable theory of Nature until sufficient principle reduction has been achieved.---As Russell said before \cite{Russell1961},
\begin{quote}
I cannot believe---and I say this with all the emphasis of which I am capable---that there can ever be any good excuse for refusing to face the evidence in favour of something unwelcome. It is not by delusion, however exalted, that mankind can prosper, but only by unswerving courage in the pursuit of truth.
\end{quote}

With good reason, Barnes calls the part of his book dealing with the unfalsifiable Eleatic theory `The Serpent' \cite{Barnes1982}, as its Devilish charm has seduced lazy and obtuse metaphysicists ever since. We find, however, in full examination of the situation, that Heraclitus' principle must be included in a truly objective deduction of the known facts; and, if neither theory is favoured when all others have been considered, we say that this is indeed the critical point---for it manages to tip the scale vertically, in favour of the Heraclitean theory. For this reason, the Eleatic doctrine has never truly become the paradigm in natural science, which seems often to find itself taking a Sceptical stance when faced with such theoretical inconsistencies as the Eleatics discovered.
\section{Accession to a Natural Epistemology\label{sec_ANE}}
In everything that has now been said against the Eleatic doctrine, we find ourselves in the company of good authority, as the main points made here were similarly stated by Aristotle at the beginning of his \textit{Physics}, as he, too, argued that the physical description of Nature must be based on natural principles \cite{Aristotle1984} (184b26-185a14):\footnote{See also Aristotle's \textit{Physics} 253a32-b6 \cite{Aristotle1984}.}
\begin{quotation}
Now to investigate whether what exists is one and motionless is not a contribution to the science of nature. For just as the geometer has nothing more to say to one who denies the principles of his science---this being a question for a different science or for one common to all---so a man investigating \textit{principles} cannot argue with one who denies their existence. For if what exists is just one, and one in the way mentioned, there is a principle no longer, since a principle must be the principle \textit{of} some thing or things.

To inquire therefore whether what exists is one in this sense would be like arguing against any other position maintained for the sake of argument (such as the Heraclitean thesis,\footnote{Presumably, as mentioned further on (185b19-25) \cite{Aristotle1984}: `But if all things are one in the sense of having the same definition, like raiment and dress, then it turns out that [Melissus and Parmenides] are maintaining the Heraclitean doctrine, for it will be the same thing to be good and to be bad, and to be good and to be not good; and so the same thing will be good and not good, and man and horse \ldots.' We can assume from this statement, that Aristotle saw no epistemological merit in Heraclitus' rant about the unity of contraries, and therefore gave no recognition, as we have above, that his method of Natural inquiry had anticipated the one that is currently being argued for.} or such a thesis as that which exists is one man) or like refuting a merely contentious argument---a description which applies to the arguments both of Melissus and of Parmenides: their premisses are false and their conclusions do not follow. Or rather the argument of Melissus is gross and offers no difficulty at all: accept one ridiculous proposition and the rest follows---a simple enough proceeding.

We, on the other hand, must take for granted that the things that exist by nature are, either all or some of them, in motion---which is indeed made plain by induction. Moreover, no one is bound to solve every kind of difficulty that may be raised, but only as many as are drawn falsely from the principles of the science: it is not our business to refute those that do not arise in this way \ldots
\end{quotation}

The everlasting impression of an idea that can never be refuted, and a Devilish seduction that has always found a way to draw in some Idealist sympathisers,---which has been thought by some Sceptics to be as much correct as it isn't,---is one way of looking at the Eleatic legacy; but a fairer assessment would be that the Eleatics had a logically self-consistent theory, however immature, and that they supported it with a good argument against Scepticism.

For, as the Eleatics reasoned, the way to truly refute a theory is not to criticise its inconsistency with traditional opinion, but to fully examine the arguments for and against it, and decide, based on that, which is the most reasonable position to maintain. This is the main argument that is disseminated in Plato's \textit{Parmenides}, and the point was in fact explicated by Parmenides in his poem, which was written in two such parts: \textit{The Way of Truth} (quoted above) and \textit{The Way of Opinion} (which is mostly lost).

Although the Sceptics also have a reasonable point, which shall be propounded here as well, that absolute Truths cannot be known with absolute certainty, and a theory's overall consistency with empirical evidence must be ultimately required, a Sceptic contributes nothing meaningful to our understanding of Nature with the paltry argument, that a theory cannot be correct if it is inconsistent with this premiss---and, in fact, often represses the search for a more correct theory by voicing doubts which are based on nothing more than a difference of opinion, rather than attempting to sort out what the crux of that difference really is. Accordingly, many academics in the sixteenth and seventeenth centuries were unwavering in their adherence to the geocentric principle.

Conversely, although we shall remain consistent with some basic principles of Scepticism, the aim of the present argument shall be more like that of the Eleatics; for, rather than allowing the inner Sceptic to merely voice obvious criticisms, the primary aim here has been to clearly \textit{understand}, through logical analysis, the precise reasons why very brilliant people allowed themselves to accept theories of Nature which are obviously incorrect; for we would be at the height of ignorance to presume that geniuses like Parmenides and Einstein did not have very good reasons for propounding the theory that time is merely an illusory aspect of a four-dimensional `Block'. And it will be clearly shown how wrong it has been to ignore the fact that this \textit{is} what their theories formally required, and go on using them to explain processes which are patently dynamical; for, by analysing both sides of the argument in the Eleatic tradition, we shall be able to determine where the true points of conflict lie, and see clearly that the principles which end up formally requiring a strictly determined, substantive block, are precisely the ones that have also led to significant mental blocks, as paradoxes have been subsequently inferred when attempts were made to describe dynamical processes by the same theories.

So we acknowledge the significant contribution of the Eleatics, who developed the method of using deductive reasoning in order to derive a logical theory that is consistent with basic principles. Once any such theory exists, it is possible to use scientific induction, to state testable hypotheses which allow us to refine our knowledge of Nature as important discoveries happen to be made. As Thomas Kuhn explained, however, this process of refinement is not continuous, as, along with the discoveries which come in line with the current paradigm, significant problems are also continually discovered, which eventually lead to paradigm shifts, when a new theory is discovered which brings that knowledge all into line. To quote the blurb from his seminal essay on \textit{The Structure of Scientific Revolutions} \cite{Kuhn1962},
\begin{quote}
Drawing his data from history, philosophy, and psychology, Kuhn argues that ``normal science'' presupposes a conceptual and instrumental framework or para-digm accepted by an entire scientific community; that the resulting mode of scientific practice inevitably evokes ``crises'' which cannot be resolved within this framework; and that science returns to normal only when the community accepts a new conceptual structure which can again govern its search for novel facts and for more refined theories.
\end{quote}

It is, in fact, reasonable to suggest that these paradigm shifts might ultimately arise from discoveries made according to the less commonly practised Heraclitean method of introspection, or principle reduction, which is a necessary element of the method of inquiry required to form a complete understanding of the workings of Nature; consequently, it would seem that only through the marriage, and continued practice, of these two essential methods,---in probing and discerning the potentially basic principles of Nature, and then deducing greater theories in agreement with the novel facts of observation, which ultimately leads to further testable hypotheses,---does the paradigm of natural philosophy truly progress---i.e., that by systemmatically probing the limits of logically deduced, generally consistent metaphysical theories by both induction and reduction, we form a clear and self-consistent idea of what Nature really is.

In truth, though, Heraclitean reduction is often neglected in practice, whenever the current paradigm is assumed as an absolute Truth. Then, it seems that the paradigm itself eventually becomes foggy, due to a common neglect of the necessary continuous introspective rationalisation, and then unclearly defined, due to a subsequent necessity to reconcile the paradoxes that ultimately result when inductive progress is not matched by the necessary reduction, or a general deduction has been neglected when parts of the theory have become too specialised and the consistency of the whole has thus become difficult to infer; and science allows itself to sumble about, half-blind, through an interval bounded at the one end by an unsure, yet comfortably assumed bias, and at the other, by the boundary that assumption places on the possibilities of theoretical advancement---for the two ends are really parts of a continuous interval, of which progress can only go so far in one direction, without seeing an equivalent amount in the other. 

Thus, the historical evidence has shown that after any paradigm has enjoyed the comfort of its respite, and long tradition has replaced original reasoning, as the agent of its acceptance---after blind stumbling has eventually allowed progress to be made in one direction only, having uncovered as many novel facts and theoretical refinements as possible, yet these are accompanied by an equivalent amount of problems---we see that the most sure way to progress when the continued practice of scientific induction alone has led to the realisation of significant problems or inconsistencies within the current paradigm, is \textit{not} to continue piling \textit{a posteriori} hypotheses and conjectures onto that paradigm at leisure, nor to propose baseless untestable new theories with naturally uninspired mathematical finesse and hopeless desparation---in both cases relying on blind luck and a remote possibility of some eventual observation, etc., for validation of such educated guesswork,---but to rationally examine the basic principles of the current paradigm, with the goal of further clarifying what they must be, and why; and possibly, with that clearer understanding, reducing them to more basic ones, and thereby emerging with a new and clearer view of Physical Reality, as well as a theory that is consistent with the old one, but which also accounts for the problems that have been realised through inductive inference from the empirical data, according to subsequent objective deduction.

To state this critical point more plainly, it is not enough to try to objectively interpret our observations, which interpretations inevitably conform to the current paradigm, even if they are thereby understood to be paradoxical, but we must also constantly, prudently reason our way to the Truth. For the ultimate goal of all this, is not merely a model deduced from the current theory which would account for the observed facts; and it is cheap, and therefore generally ineffectual, to create a whole new unverifiable theory out of nothing but imprudent philosophical speculation (even if that should be mathematical in nature) that has no prior connection, common sense-based or otherwise, to a realistic understanding of nature. The goal is in fact to understand `Why?' And this question cannot be answered without a clearly reduced metaphysical theory based on concrete natural principles, from which such a model---i.e., one that would account for all the observed facts---should naturally emerge, according to those prior natural causes. 

The age-old hitch in actually achieving this goal, as we know, is the inability or unwillingness of subsequent practitioners, to refine their basic understanding once a problem with the accepted physical theory has been realised, as they would still rather work within the bounds that are set by the paradigm it created; or, as Einstein put it in his 1916 tribute to Ernst Mach \cite{Einstein1916},
\begin{quotation}
%Begriffe, welche sich bei der Ordnung der Dinge als nuetzlich erwiesen haben, erlangen ueber uns leicht eine solche Autoritaet, dass wir ihres irdischen Ursprungs vergessen und sie als unabaenderliche Gegebenheiten hinnehmen. Sie werden dann zu "Denknotwendigkeiten", "Gegebenen a priori", usw. gestempelt. Der Weg des wissenschaftlichen Fortschrittes wird durch solche Irrtuemer oft fuer lange Zeit ungangbar gemacht. Es ist deshalb durchaus keine muessige Spielerei, wenn wir darin geuebt werden, die laengst gelaeufigen Begriffe zu analysieren und zu zeigen, von welchen Umstaenden ihre Berechtigung und Brauchbarkeit abhaengt, wie sie im einzelnen aus den Gegebenheiten der Erfahrung herausgewachsen sind. Dadurch wird ihre allzu grosse Autoritaet gebrochen. Sie werden entfernt, wenn sie sich nicht ordentlich legitimieren koennen, korrigiert, wenn ihre Zuordnung zu den gegebenen Dingen allzu nachlaessig war, durch andere ersetzt, wenn sieh ein neues System aufstellen laesst, das wir aus irgendwelchen Gruenden vorziehen.

%Derartige Analysen erscheinen dem Fachwissenschaftler, dessen Blick mehr auf das Einzelne gerichtet ist, meist ueberfluessig, gespreizt, zuweilen gar laecherlich. Die Situation aendert sich aber, wenn einer der gewohnheitsmaessig benutzten Begriffe durch einen schaerferen ersetzt werden soll, weil es die Entwicklung der betreffenden Wissenschaft erheischt. Dann erheben diejenigen, welche den eigenen Begriffen gegenueber nicht reinlich verfahren sind energischen Protest und klagen ueber revolutionaere Bedrohung der heiligsten Gueter. In dies Geschrei mischen sich dann die Stimmen derjenigen Philosophen welche jenen Begriff nicht entbehren zu koennen glauben, weil sie ihn in ihr Schatzkaestlein des "Absoluten" des "a priori" oder kurz derart eingereiht hatten, dass sie dessen prinzipielle Unabaenderlichkeit proklamiert hatten.
\ldots. Concepts which have proven useful in ordering things, easily attain such an authority over us that we forget their Earthly origins and accept them as unalterable facts. They are then branded as `necessities of thought', `a priori givens', etc. The path of scientific advance is often made impassable for a long time through such errors. It is therefore by no means an idle trifling, if we become practiced in analysing the long-familiar concepts, and to show upon which circumstances their justification and applicability depend, as they have grown up, individually, from the facts of experience. For through this, their all-too-great Authority will be broken. They will be removed, if they cannot be properly legitimated, corrected, if their correlation to given things was far too careless, or replaced by others, if we see a new system that can be established, that we prefer for whatever reasons.

This type of analysis appears to the scholars, whose gaze is directed more at the particulars, most superfluous, splayed, and at times even ridiculous. The situation changes, however, when one of the habitually used concepts should be replaced by a sharper one, because the development of the science in question demanded it. Then, those who are faced with the fact that their own concepts do not proceed cleanly raise energetic protest and complain of revolutionary threats to their most sacred possessions. In this cry, then, mix the voices of those philosophers who believe those concepts cannot be done without, because they had them in their little treasure chest of the `absolute', the `a priori', or classified in just such a way that they had proclaimed the principle of immutability.
\end{quotation}

And so it is regrettable that physicists have not since made greater strides to incorporate such investigations of the ontological side of physical theory into regular practice, along with the study and development of appropriate epistemological methods, and thus work to close the gap between philosophy and science as physical theory has developed, rather than leaving the philosophical part of theoretical development to be investigated only by the philosophers whose opinions they're apt to ignore; for, as Einstein also wrote in this article, from the height of scientific achievement \cite{Einstein1916}, 
\begin{quotation}
%Wie kommt aber ein ordentlich begabter Naturforscher ueberhaupt dazu, sich um Erkenntnistheorie zu kuemmern? Gibt es nicht in seinem Fache wertvollere Arbeit? So hoere ich manche meiner Fachgenossen hierauf sagen, oder spuere bei noch viel mehr, dass sie so fuehlen. Diese Gesinnung kann ich nicht teilen. Wenn ich an die tuechtigsten Studenten denke, die mir mein Lehren begegnet sind, d. h. an solche, die sich durch Selbstaendigkeit des Urteils, nicht nur durch blosse Behendigkeit auszeichneten, so konstatiere ich bei ihnen, dass sie sich lebhaft um Erkenntnistheorie kuemmerten. Gerne gegangen sie Diskussionen ueber die Ziele und Methoden der Wissenschaften und zeigten durch Hartnaeckigkeit im Verfechten ihrer Ansichten unzweideutig, dass ihnen der Gegenstand wichtig erschien. Dies ist fuerwahr nicht zu verwundern.

%Wenn ich mich nicht aus aeusseren Gründen, wie Gelderwerb, Ehrgeiz und auch nicht oder wenigstens nicht ausschliesslich des sportlichen Vergnuegens, der Lust am Gehirn-Turnen wegen einer Wissenschaft zuwende, so muss mich als Juenger dieser Wissenschaft die Frage brennend interessieren: Was fuer ein Ziel will und kann die Wissenschaft erreichen, der ich mich hingebe? Inwiefern sind deren allgemeine Ergebnisse "wahr"? Was ist wesentlich, was beruht nur auf Zufaelligkeiten der Entwicklung?

But how does it happen anyway, that a properly endowed natural scientist comes to concern himself with epistemology? Is there not more valuable work in his own trade? I hear some of my colleagues saying this, or sense from many more that they feel this way. I can not share this attitude. When I think about the ablest students I have encountered in my teaching, viz. those who distinguish themselves through independence of judgment, and not through sheer agility only, so I state of them that they had a lively interest in epistemology. They gladly entered into discussions about the aims and methods of the sciences, and showed unequivocally, through persistence in advocating their views, that the subject seemed important to them. In truth, this is not surprising.

If I am not ambitious for external reasons, such as making money, and also not, or at least not exclusively, the sporting pleasure, or delight in brain-gymnastics due to a scientific turning, then, as a disciple of this science, I must have a burning interest in the question: What possible goal does the science want to reach, to which I dedicate myself? To what extent are its general results `true'? What is essential, which is based solely on accidents of development?
\end{quotation}

The significance of this last question warrants some elaboration, so that its meaning should not be missed: in common scientific practice, we inductively infer certain properties of Nature according to empirical observation, and scientific theory grows by piecing together all of these primary inferences. But it is only by stepping back from the specifics of the theory and considering the whole of scientific and philosophic knowledge objectively, while examining the circumstances upon which the justification and applicability of the long-familiar concepts depend, that we can rationally reduce those `accidents of development' to principles which are closer to essential properties of Physical Reality. 

Accordingly, the most novel of observations, such as the discovery that the cosmic expansion is now accelerating, should not be merely fitted into the scientific paradigm that exists at the time of discovery; but reason, in conjunction with mathematics, should be applied to such problems in order to account for their consistency with the whole set of knowledge.

Now, Einstein's convictions about the essential importance of being ever-conscious of the epistemological methods one uses in conducting scientific investigation was no mere passing fancy either; as he stated unequivocally in 1949 \cite{Einstein1951},
\begin{quotation}
The reciprocal relationship of epistemology and science is of noteworthy kind. They are dependent upon each other. Epistemology without contact with science becomes an empty scheme. Science without epistemology is---insofar as it is thinkable at all---primitive and muddled. 
\end{quotation}
The reason for this, is simply that a physical theory which aims to accurately understand and describe Physical Reality cannot be complete without a corresponding clear notion (which is often achieved through elucidation) of its epistemological intent, without which it often loses sight of its purpose, its meaning, and its logical coherence. But in practice, the requirement of such constant contact with a coherent epistemological method is no trivial task to consciously maintain; and, as Einstein went on to explain after the above statement, there is as yet no strict system general enough to properly account for all the possibilities that arise in the process of fundamental inquiry \cite{Einstein1951}:
\begin{quote}
However, no sooner has the epistemologist, who is seeking a clear system, fought his way through to such a system, than he is inclined to interpret the thought-content of science in the sense of his system and to reject whatever does not fit into his system. The scientist, however, cannot afford to carry his striving for epistemological systematic that far. He accepts gratefully the epistemological conceptual analysis; but the external conditions, which are set for him by the facts of experience, do not permit him to let himself be too much restricted in the construction of his conceptual world by the adherence to an epistemological system. He therefore must appear to the systematic epistemologist as a type of unscrupulous opportunist: he appears as \textit{realist} insofar as he seeks to describe a world independent of the facts of perception; as \textit{idealist} insofar as he looks upon the concepts and theories as the free inventions of the human spirit (not logically derivable from what is empirically given); as \textit{positivist} insofar as he considers his concepts and theories justified \textit{only} to the extent to which they furnish a logical representation of relations among sensory experiences. He may even appear as \textit{Platonsit} or \textit{Pythagorean} insofar as he considers the viewpoint of logical simplicity as an indispensable and effective tool of his research.
\end{quote}

The goal of such investigations is indeed a formal analytical theory which is built on basic principles or axioms, using them to develop a mathematical formalism that would describe the empirical data; and the theoretical physicist must therefore be more than one endowed with a great arsenal of mathematical techniques:---he must be an analyst in the full sense, who rationally scrutinises his theory and further develops its possibilities through logical and mathematical deduction in order to naturally `save the phenomena'; he must be a rational (foundationalist)-empirical (idealist) analyst who therefore accepts the fact that he can never presume the Correctness of his theory, which must necessarily be built upon axioms, as well as a naturalist who seeks a logically coherent, and naturally \textit{seamless} connection between such axioms and all that can be discerned from Nature.

In his `obituary' \cite{Einstein1951}, Einstein wrote vaguely (sc., with `meager precision') of how this should be done, when he discussed the two `points of view according to which it is possible to criticize physical theories at all';---the one of which concerns itself with the `premises of the theory itself, with what may briefly but vaguely be characterized as the ``naturalness'' or ``logical simplicity'' of the premises', i.e.~`with the ``inner perfection'' of the theory';---and the other, primary point, which `refers to the ``external confirmation{''}', is `concerned with the confirmation of the theoretical foundation by the available empirical facts', he stated precisely as the requirement `that the theory must not contradict empirical facts';---and, furthermore, that this strict adherence to observation must not be achieved through `the adaptation of the theory to the facts by means of artificial additional assumptions'; for, in this way, by continually stressing both of these critical limits, and facilitating their seamless connection through logical deduction which is often realised through precise mathematical expression, `we are confining ourselves to such theories whose object is the \textit{totality} of all physical appearances'.

Therefore, the end goal of fundamental inquiry is to achieve a basic theoretical expression for our Cosmic Existence, which lends itself naturally to a description of all the observed phenomena that does not require unqualified presumption, or \textit{ad hoc} definition, of any characteristics of Nature, such as action-at-a-distance, empirically given `fundamental' constants, etc; for `there are no \textit{arbitrary} constants of this kind; that is to say, nature is so constituted that it is possible logically to lay down such strongly determined laws that within these laws only rationally completely determined constants occur (not constants, therefore, whose numerical value could be changed without destroying the theory)' \cite{Einstein1951}.

A final remark must be made before we can move beyond this discussion of Einstein's epistemological methods:---it may also be that the scientist, and even one with the noblest intentions, should fall into the trap to which Einstein condemned systematic epistemologists in \cite{Einstein1951}; for when his physical theory is finally complete, he may have become too immersed in it to be able to see how to further reduce his principles, expand the theory, and therefore see how it can be brought into greater agreement with sense-perceptions; and this indeed appears to have occurred for both Einstein and the Eleatics. With Einstein, the reason for this seems most likely to have been a caveat that warrants mention in regard to the goal of logical simplicity or logical economy of principles, which he held to be the mark of a great theory, and therefore took as his aim. The problem with this, as I see it, is that when one actively seeks out principles that would satisfy this requirement, and does not require above all the `naturalness' or `logical seamlessness' of the theory, along with the intuitive clarity that these afford, one may be led into abstraction by neglecting to properly account for certain facts of observation in this way, as Einstein was. Along with that which was noted at the end of Chapter~\ref{chap_MCCC}, this point seems to stand out, e.g., in a passage from `The Problem of Space, Ether, and the Field in Physics' (1934) \cite{Einstein1954}:
\begin{quotation}
The theory of relativity is a fine example of the fundamental character of the modern development of theoretical science. The initial hypotheses become steadily more abstract and remote from experience. On the other hand, it gets nearer to the grand aim of all science, which is to cover the greatest possible number of empirical facts by logical deduction from the smallest possible number of hypotheses or axioms. Meanwhile, the train of thought leading from the axioms to the empirical facts or verifiable consequences gets steadily longer and more subtle. The theoretical scientist is compelled in an increasing degree to be guided by purely mathematical, formal considerations in search for a theory, because the physical experience of the experimenter cannot lead him up to the regions of highest abstraction. The predominantly inductive methods appropriate to the youth of science are giving place to tentative deduction. Such a theoretical structure needs to be very thoroughly elaborated before it can lead to conclusions which can be compared with experience. Here, too, the observed fact is undoubtedly the supreme arbiter; but it cannot pronounce sentence until the wide chasm separating the axioms from their verifiable consequences has been bridged by much intense, hard thinking. The theorist has to set about this Herculean task fully aware that his efforts may only be destined to prepare the death blow to his theory. The theorist who undertakes such a labor should not be carped at as ``fanciful''; on the contrary, he should be granted the right to give free reign to his fancy, for there is no other way to the goal. His is no idle daydreaming, but a search for the logically simplest possibilities and their consequences. This plea was needed in order to make the listener or reader more inclined to follow the ensuing train of ideas with attention; it is the line of thought which has led from the special to the general theory of relativity and thence to its latest offshoot, the unified field theory.
\end{quotation}

Now, I do not wish for my point to be misunderstood here: I agree fairly well with everything in this statement. For logical economy can, and indeed should, be held in sight when searching for a basic theory of Nature; but it cannot be the \textit{ultimate} goal of the metaphysicist (for that is what one is, if one begins to rationally scrutinise the meaning of physical theory and its connection to Physical Reality, rather than merely conducting experiments based on the prior theorisation of others, with the goal of making inductive inferences to advance that theory) to derive the most economical principles with the farthest reaching consequences, if that would come with the price of abstraction due to a neglect of any reproducible empirical datum---and especially not one that was accounted for in the older theory;---for that is no proper reduction. 

And the fact is, that the Newtonian concept of an absolute time, though too carelessly formulated in the mathematical theory, is in accord with Heraclitus' principle, and thus with a rationalisation of common sense-perception that is no less pertinant to our theoretical description of Physical Reality than any physical interpretation of an observation that might be made, e.g., at the Large Hadron Collider; therefore, the Newtonian theory possesses an intuitive clarity which the theory of relativity, according to the common interpretation that leads to a block universe description in no less straightforward a manner than the Eleatic deduction, does not. As Einstein wrote in his autobiography \cite{Einstein1951},
\begin{quote}
Newton, forgive me; you found the only way which, in your age, was just about possible for a man of highest thought- and creative power. The concepts, which you created, are even today still guiding our thinking in physics, although we know that they will have to be replaced by others farther removed from the sphere of immediate experience, if we aim at a profounder understanding of relationships.
\end{quote}

The notion that certain principles can be simply abandoned for others in order to achieve the `grand aim of all science', is a recurring theme in Einstein's epistemology, and a mark against it. For in order to properly reduce the concepts, hypotheses, or axioms of any theory that does hold the weight of Authority, as it had previously been thought to provide a good description of Nature, \textit{only} the second option that Einstein stated in his tribute to Mach can be allowed; viz., they must be corrected, by showing that their correlation to the given things was far too careless. For then it is reasonable to expect that logical economy shall be an accidental property of the resulting theory.

The epistemological method discussed here, is in fact the method that has been practised, to some degree or other, by all of the greatest thinkers in history who have made the most important advances in our understanding of Nature. In ancient times, we see its mark in Democritus' search for Truth, the worth of which he placed above the kingdom of the Persians, and in the general completeness of his theory, which he derived by rational thought based on primary natural observations; and in Aristotle, who not only meticulously deduced a physical theory based on such prior natural principles, that was as epistemologically complete as any has been since, but who also exhaustively formalised the requisite methods of logical inquiry; and we know this to have been a central aspect of Epicurean epistemology, as they devoted their study to the discovery of principal doctrines that could serve both, the greater understanding of the whole, and, subsequently, more specialised inquiry. We see this, too, in the works of the great natural philosophers of the Scientific Revolution---those apostates of `Aristotelianism' and seekers of Truth whose tremendous achievement in bringing philosophy out of its darkest Age was epitomised by the efforts of Kepler and Galileo, and finally culminated in the great Newtonian synthesis; and we also find its mark in the list of brilliant minds that followed, from which the names, Darwin, Maxwell, and Einstein, stand out above all others as the most successful inferers of natural principles, who managed to deduce, from these and the greater set of physical observations, general theories of Nature that would further revolutionise our understanding of Its workings.

For along with Einstein, whose methodology we have now considered in some detail, it is equally true that Charles Darwin managed to separate himself from the common theory of his time, which required that every species of life on Earth must remain essentially unchanged subsequent to its Creation, and recognised that the awesome complexity of the great system of species could be more reasonably accounted for by a theory in which all the species have evolved from a few, or even one progenitor, according to the principle of natural selection. And, in the preface to his great \textit{Treatise on Electricity and Magnetism}, James Clerk Maxwell laid out the method by which he formulated the electromagnetic theory, which exemplifies many aspects of the epistemological method that has now been discussed; for there he began with a statement of how he would deduce, from all the data, the general mathematical theory that would describe elecrtomagnetic phenomena with respect to observable quantities, with the aim of reducing electromagnetism to one interconnected, self-consistent dynamical theory; then, more specifically, he wrote \cite{Maxwell1892},
\begin{quotation}
I have \ldots thought that a treatise would be useful which should have for its principal object to take up the whole subject in a methodical manner, and which should also indicate how each part of the subject is brought within the reach of methods of verification by actual measurement.

The general complexion of the treatise differs considerably from that of several excellent electrical works, published, most of them, in Germany, and it may appear that scant justice is done to the speculations of several eminent electricians and mathematicians. One reason of this is that before I began the study of electricity I resolved to read no mathematics on the subject till I had first read through Faraday's \textit{Experimental Researches in Electricity}. I was aware that there was supposed to be a difference between Faraday's way of conceiving phenomena and that of the mathematicians, so that neither he nor they were satisfied with each other's language. I had also the conviction that this discrepancy did not arise from either party being wrong. \ldots

As I proceeded with the study of Faraday, I perceived that his method of conceiving the phenomena was also a mathematical one, though not exhibited in the conventional form of mathematical symbols. I also found that these methods were capable of being expressed in the ordinary mathematical forms, and thus compared with those of the professed mathematicians.

\ldots

When I had translated what I considered to be Faraday's ideas into a mathematical form, I found that in general the results of the two methods coincided, so that the same phenomena were accounted for, and the same laws of action deduced by both methods, but that Faraday's methods resembled those in which we begin with the whole and arrive at the parts by analysis, while the ordinary mathematical methods were founded on the principle of beginning with the parts and building up the whole by synthesis.

\ldots

\ldots. The great success which these eminent men [sc.~the German electricians and mathematicians] have attained in the application of mathematics to electrical phenomena, gives, as is natural, additional weight to their theoretical speculations, so that those who, as students of electricity, turn to them as the greatest authorities in mathematical electricity, would probably imbibe, along with their mathematical methods, their physical hypotheses.

These physical hypotheses, however, are entirely alien from the way of looking at things which I adopt, and one object which I have in view is that some of those who wish to study electricity may, by reading this treatise, come to see that there is another way of treating the subject, which is no less fitted to explain the phenomena, and which, though in some parts it may appear less definite, corresponds, as I think, more faithfully with our actual knowledge, both in what it affirms and in what it leaves undecided.
\end{quotation} 

Thus, Maxwell recognised the philosophical brilliance of Michael Faraday, who began by considering the whole system of observed electromagnetic phenomena, and inferred a general framework for electromagnetic field theory which is consistent overall. He then applied himself to understanding the phenomena in this way, and then used mathematical expertise to analytically deduce the corresponding theory.

This is basically the same approach that has been taken more recently by David Bohm and Basil Hiley, in \textit{The Undivided Universe}. As Hiley wrote in the Preface \cite{Bohm1993},
\begin{quote}
I hope this book will be a fitting testimony to this very radical and original thinker [Bohm] who rejected the view of conventional quantum mechanics, not for ideological reasons, but because it did not provide a coherent overall view of nature, a feature that David felt an essential ingredient of any physical theory. It was like Escher's ``The Waterfall'', a fascinating picture in which region by region appeared to be carefully constructed and consistent, but when one stepped back to perceive the whole, a contradiction was there for all to see. Indeed the most radical view to emerge from our deliberations was the concept of wholeness, a notion in which a system formed a totality whose overall behaviour was richer than could be obtained from the sum of its parts.
\end{quote}
As well, they say later on, regarding the conventional interpretation of quantum mechanics, that `It is not that we want to dismiss the epistemological approach totally, but rather that we feel that both approaches have to be pursued and perhaps ultimately brought into relationship.' For, as they said in the introduction, one of the main advantages of their interpretation is that
\begin{quote}
it provides an intuitive grasp of the whole process. This makes the theory much more intelligible than one that is restricted to mathematical equations and statistical rules for using these equations to determine the probable outcomes of experiments. Even though many physicists feel that making such calculations is basically what physics is all about, it is our view that the intuitive and imaginative side which makes the whole theory intelligible is as important in the long run as is the side of mathematical calculation.
\end{quote}

One final quotation from the book should help to round off their argument, and show that it is pretty much in line with what has been written above:
\begin{quotation}
\ldots Heisenberg \ldots said very explicitly that the essential truth was in the mathematics. This view has become the common one among most of the modern theoretical physicists who now regard the equations as providing their most immediate contact with nature (the experiments only confirming or refuting the correctness of this contact). On the other hand, in the past we began talking about our concept of physical reality and used the equations to talk about them.

One of the reasons that our interpretation of the quantum theory may not have been so well understood is perhaps that it is not commonly realised that we have a rather different attitude to the mathematics. Our approach is not simply a return to the notion that the mathematics merely enables us to talk about the physical concepts more precisely. Rather we feel that these two kinds of concepts represent extremes and that it is necessary to be in a process of thinking that moves between these extremes in such a way that they complement each other.

We fully accept that progress can be made from the mathematical side alone, but we feel that to stick indefinitely to this approach is too limited. \ldots we do not regard \ldots physical concepts as mere imaginative displays of the meaning of our equations. Instead, as we shall see in this chapter, we can then move to the other side of this process of thought and take our physical concepts as a guide for the development of new equations. In this way we are engaged in at least the beginning of a kind of thinking in which the clue for a creative new approach may come from either side and may flow indefinitely back and forth between the two sides.
\end{quotation}

Actually, the most significant difference between Bohm and Hiley's ontological outlook and the one taken here, is that, although they defend the concept of an ultimate Reality, they argue by inductive inference that it is `unlimited and unknown', whereas we shall take the alternate view-point, that there are well-defined, and therefore potentially definable, Truths of Nature, which could be stated as mathematical principles of natural philosophy. But although such principles are essential to any complete physical theory, regardless of how they might be connected to an ultimate Reality, it is also recognised, as a basic property of logical inquiry, that we can never be absolutely certain whether the principles or axioms we've assumed are actually True.

Perhaps the clearest systematic application of the method discussed here, which must be used in order to determine not merely a physical theory that would describe the observed phenomena, but one that works from a clear realisation of basic natural principles to deduce a complete and self-consistent theory, which is nevertheless fully aware of its weaknesses, is found in Isaac Newton's \textit{Philosophi{\ae} Naturalis Principia Mathematica}. Newton begins the \textit{Principia}, following some definitions of less commonly known physical quantities, such as mass, momentum, inertia, mechanical forces, etc., and a description of what he would precisely mean by the better-known concepts of time, space, place, and motion, with a statement of three basic mathematical principles, his famous Laws of motion, before going on, in the first two books, to use them in order to describe in detail the kinematical motions of bodies. Then, he begins the third book with the statement \cite{Newton1966},
\begin{quotation}
In the preceding books I have laid down the principles of philosophy; principles not philosophical but mathematical: such, namely, as we may build our reasonings upon in philosophical inquiries. These principles are the laws and conditions of certain motions, and powers or forces, which chiefly have respect to philosophy; but, lest they should have appeared of themselves dry and barren, I have illustrated them here and there with some philosophical scholiums, giving an account of such things as are of more general nature, and which philosophy seems chiefly to be founded on; such as the density and the resistance of bodies, spaces void of all bodies, and the motion of light and sounds. It remains that, from the same principles, I now demonstrate the frame of the System of the World. \ldots
\end{quotation} 
He eventually moves on to write down four Rules of reasoning in philosophy:
\begin{itemize}
\item[I]{`We are to admit no more causes of natural things than such as are both true and sufficient to explain their appearances.'}
\item[II]{`Therefore to the same natural effects we must, as far as possible, assign the same causes.'}
\item[III]{`The qualities of bodies, which admit neither intensification nor remission of degrees, and which are found to belong to all bodies within reach of our experiments, are to be esteemed the universal qualities of all bodies whatsoever.'}
\item[IV]{`In experimental philosophy we are to look upon propositions inferred by general induction from phenomena as accurately or very nearly true, notwithstanding any contrary hypotheses that may be imagined, till such time as other phenomena occur, by which they may either be made more accurate, or liable to exceptions.'}
\end{itemize}
Then, after a lengthy analysis involving his Law of universal gravitation, he concludes with a general scholium, where he criticises Ren\'{e} Descartes' vortex theory, and then adds some unjustified speculation which I'll have to come back to, before concluding with,
\begin{quotation}
\ldots hitherto I have not been able to discover the cause of those properties of gravity from phenomena, and I frame no hypotheses; for whatever is not deduced from the phenomena is to be called an hypothesis; and hypotheses, whether metaphysical or physical, whether of occult qualities or mechanical, have no place in experimenal philosophy. In this philosophy particular propositions are inferred from the phenomena, and afterwards rendered general by induction. Thus it was that the impenetrability, the mobility, and the impulsive force of bodies, and the laws of motion and of gravitation, were discovered. And to us it is enough that gravity does really exist, and act according to the laws which we have explained, and abundantly serves to account for all the motions of the celestial bodies, and of the sea.

And now we might add something concerning a certain most subtle spirit which pervades and lies hid in all gross bodies; by the force and action of which spirit the particles of bodies attract one another at near distances, and cohere, if contiguous; and electric bodies operate to greater distances, as well repelling as attracting the neighboring corpuscles; and light is emitted, reflected, refracted, inflected, and heats bodies; and all sensation is excited, and the members of animal bodies move at the command of the will, namely, by vibrations of this spirit, mutually propagated along the solid filaments of the nerves, from the outward organs of sense to the brain, and from the brain into the muscles. But these are things that cannot be explained in few words, nor are we furnished with that sufficiency of experiments which is required to an accurate determination and demonstration of the laws by which this electric and elastic spirit operates.
\end{quotation}

If we consider this final statement along with his Rules of philosophy, it becomes clear that Newton's theory also contained the necessary statement of a greater epistemological purpose, to continue with experimental research until such time as sufficient data existed which would indicate some direction for a subsequent reduction of principles more fundamental than the laws he was capable of inferring in his time; for he admittedly did not know the cause of those laws, and rightly framed no hypotheses to account for them. But he would have done as well to have adhered strictly to this professed stance, and not allow himself to be victimised by dogmatic prejudice into a direct contradiction; as, just prior to this concluding statement, he \textit{had} framed a related hypothesis, which, above, I've referred to as `unjustified speculation'---i.e., that the `most beautiful system of the sun, planets, and comets, could only proceed from the counsel and dominion of an intelligent and powerful Being.' For, just as the law of universal gravitation, and the equivalence of gravitational and inertial mass, are better understood as they are explained by general relativity theory: so, too, within the modern theory of cosmic evolution, relativistic cosmology is found to be sufficient to give rise naturally to the great system of the \textit{Universe}---the clusters of galaxies and all their subsystems,---as having formed dynamically since the Big Bang---which is one of the best examples of the achievement of Einstein's theory over Newton's.

In contrast to this appeal to the `counsel and dominion' of God, in order to explain the formation of our Solar System, even though he would `frame no hypotheses' about his law of universal gravitation, it is interesting to note the stance taken by Kepler in the introduction to his \textit{Astronomia Nova} (where he derived the first two of his laws), in 1609 \cite{Koestler1959}:
\begin{quotation}
So much for the authority of Holy Scripture. Now as regards the opinions of the saints about these matters of nature, I answer in one word, that in theology the weight of Authority, but in philosophy the weight of Reason alone is valid. Therefore a saint was Lactantius, who denied the earth's rotundity; a saint was Augustine, who admitted the rotundity, but denied that antipodes exist. Sacred is the Holy Office of our day, which admits the smallness of the earth but denies its motion: but to me more sacred than all these is Truth, when I, with all respect for the doctors of the Church, demonstrate from philosophy that the earth is round, curcumhabited by antipodes, of a most insignificant smallness, and a swift wanderer of the stars.
\end{quotation}
Or, we might consider, from Galileo's letter of acknowledgement to Kepler, upon receiving a copy of his \textit{Mysterium Cosmographicum} in 1597 (a rare courtesy that was paid in that direction) \cite{Koestler1959},
\begin{quotation}
\ldots I indeed congratulate myself on having an associate in the study of Truth who is a friend of Truth. For it is a misery that so few exist who pursue the Truth and do not pervert philosophical reason. \ldots
\end{quotation}

To be sure, though: Kepler's contributions to modern science were not the results of strict adherence to this dictum, from which he frequently departed \cite{Koestler1959}; whereas Newton's speculation about the formation of the Solar System was rather an exception to the rule.\footnote{As one of Kepler's biographers has remarked, `Not the least achievement of Newton was to spot the Three Laws in Kepler's writings, hidden away as they were like forget-me-nots in a tropical flowerbed' \cite{Koestler1959}.} And aside from that, it should be noted---say, according to his Rules of reasoning in philosophy, which we see no cause to go against,---that Newton, with good reason, also supported the assumption that there \textit{is} some basic physical Truth which holds always, and that there are, therefore, basic Natural Principles which are properties of that Truth; i.e., that there are basic physical Truths which hold consistently in Reality, throughout space and time. This, we'll refer to as the principle of Natural Order. Accordingly, we define a \textit{natural principle} as a (theoretic) principle derived through the inductive process of Heraclitean introspection, or insightful rationalisation, which might later be reduced to a principle which is logically more applicable to the scope of scientific understanding. We then deduce our theories based on the further assumption that these principles \textit{could} be Natural Truths, and, if the resulting theory is consistent with all known natural principles, we use the scientific method to examine it for overall self-consistency, and to further broaden its scope. 

However, if the theory is not consistent with all known natural principles, we recognise a significant philosophical problem, like any other we might discover through scientific investigation. We say that any theory, however powerful it may be, shall not be considered even provisionally correct unless it is consistent with every natural principle and fact within the current field of knowledge, and then only as long as new problems are not discerned; i.e.---borrowing from Leibniz' principle of sufficient reason, and from metaphysical naturalism in general, we assume an epistemological \textit{principle of natural cause} (which is similarly expressed in Newton's Rules of philosophical reasoning),---that \textit{there must be a basic reason for every fact of observation which is ultimately consistent with everything else we can know of Nature, and is sufficient to explain why it should be so, and not otherwise}.

Thus, a challenge comes down to us: that, just as Maxwell did, we must abandon that inherent stubborn and na\"{i}ve optimism of prejudiced thinking, that allows us to be convinced of the absolute Truth of any basic metaphysical truth of any current scientific paradigm, even in any respect we are \textit{certain} we know it to be true, which is a constant crutch, and a general oppressor of the search for more basic knowledge,---and replace this with a rational optimism and clear acknowledgement that our theories \textit{could} be True, as that may be warranted, but that theoretical truth must be always carefully scrutinised as our theories continue to develop; i.e., that along with making inductive interpretations of scientific observations, we must also work by continually reasoning our way towards truths which are consistent with increasingly large sets of empirical data. For even when a scientific theory is known to be entirely self-consistent, and is thus to be considered theoretically true, as it is in accord with all current observations and natural principles, we cannot ever treat it as the absolute Truth, because it will remain the product of logical analysis. This was wonderfully expressed by Einstein in the opening of his centenary commemoration of Maxwell, when he said \cite{Thomson1931},
\begin{quotation}
The belief in an external world independent of the percipient subject is the foundation of all science. But since our sense-perceptions inform us only indirectly of this external world, or Physical Reality, it is only by speculation that it can become comprehensible to us. From this it follows that our conceptions of Physical Reality can never be definitive; we must always be ready to alter them, to alter, that is, the axiomatic basis of physics, in order to take account of the facts of perception with the greatest possible logical completeness. \ldots
\end{quotation}

But the reassurance of such provisional epistemological completeness, as, e.g., the Newtonian theory had once possessed, as far as it was capable of accurately describing various idealised phenomena according to four basic Laws, is currently not afforded by either of the two great theories of modern physics: quantum theory, which of course does work from some principles, and which has been extremely successful, possesses no basic natural principles, and is therefore subject to multiple possible interpretations, and many paths of reduction have thus been proposed---i.e., it is not a working \textit{metaphysical} theory;---general relativity theory, on the other hand, which was derived from natural principles, and has been equally as successful, is problematic in many ways. For these reasons, however, the methods of natural epistemology should experience far more difficulty in trying to address quantum theory directly, because we really don't have any certain indication as to its actual connection to Physical Reality, than in assessing general relativity theory, the problems of which arose because it incessantly \textit{wants} to be something other than what Einstein tried to make it, which is basically unnatural---i.e., although it is a deductive theory, it is so \textit{only} in the same fashion as the Eleatic theory, as it was never properly reconciled with Heraclitus' principle of flux, but holds fast to a seemingly natural principle that will likewise have to be abandoned once we've examined the requirements of Real dynamics more carefully. For as it currently stands, relativity is not a \textit{natural} metaphysical theory, because it is basically inconsistent with the principle of natural cause.

This problem has always been known, and many attempts over the years have tried to render relativity theory more natural than it is described in Einstein's (or, rather, Minkowski's) interpretation, by those who have recognised how unnatural his theory really was made to be. Principal among these, I think, are the works of Weyl \cite{Weyl1924}, Mili\v{c} \v{C}apek \cite{Capek1961}, and Stein \cite{Stein1968,Stein1991}, as well as the geometrodynamics theory, described in \cite{Witten1962,Arnowitt2008} and the works of John Wheeler throughout the 1960s---see \cite{Graves1971} for a comprehensive account of this theory from 1971. However, the general problem with these descriptions, I believe, is that although all fully realised the natural contradiction in Einstein's initial interpretation, none of these found a way to clearly discern what was really always at the heart of the problem---although I think Weyl, who, in \cite{Weyl1924}, did examine the natural principles of the theory in the way just described, came very close, and probably had roughly the correct intuitive understanding, though it had not been sufficiently explicit.

Of course, if the different interpretation of relativity theory we are now set to describe, is to be consistent with the current paradigm---in line with what it wants to be, but has not actually achieved, as I've just said,---it must not merely be self-consistent, but consistent with the other theories that contribute to that paradigm; specifically, it would be a bad start to begin by being inconsistent with quantum theory, as relativity theory is commonly thought to be. So, as the problem currently under investigation is the Nature of spacetime, and as we currently do not have any \textit{good} reason to believe otherwise, we shall provisionally avoid the problem of not knowing what quantum theory Really is, and therefore work towards developing a theory that may, for all we currently know, be potentially True, and thus, potentially reconcilable with a more complete understanding of quantum theory, in the following way:---we propose, as basic natural principles, the subsistence and persistence of Reality; furthermore, we will assume of that Reality, that matter and space(time) coexist in independently ontologically Real states, that space(time) is \textit{complete}, in the formal sense of metric spaces, and that the quantum field incidentally describes material interactions \textit{only}; i.e., that space(time) is \textit{not} quantised, though particle interactions may be described as such. Accordingly, by considering only the implicit Nature of spacetime to begin with, we shall not find a need to directly address quantum mechanics from the outset.\footnote{It is worth noting, that in their attempt to construct a quantum theory upon very similar grounds, Bohm and Hiley wrote, `If this theory is intended to apply cosmologically, it is evidently necessary that we should not, from the outset, assume essential elements that are not capable of being included in the theory' \cite{Bohm1993}.} 

The consistency of these assumptions shall not be directly addressed until much later in the analysis, but they are to be considered, for now, as forming a basic working principle from which we will eventually derive a theory that can be shown to be truly consistent in so many ways with the novel facts, theoretical refinements, and \textit{implicit} resolution of conceptual problems of modern physics. So now, keeping in mind both this principle and the greater epistemology to which it has been decided that all natural metaphysical theories must ultimately be objectively assessed, we shall return to our discussion of the historical progression of metaphysical theories concerning the Nature of Time.
\section{On Ancient Atomism\label{sec_OAA}}
We've seen that in his time, Parmenides had renounced Time, and thus effectively put an end to \textit{purpose} with his logical argument, if a hole could not be found in it. For if all is One, then the will that drives our efforts is purely artificial or illusory. But although the Eleatic doctrine would maintain a healthy tradition to this day, that vital snag, as I've mentioned above, was not long in coming. Not surprisingly, though, it came from Leucippus, who many accounts tell was to have come from Elea, where he had been a student of either Zeno or Melissus, who are Parmenides' two most famous followers. He is reported to have left Elea and moved to Abdera, where he developed the first atomic theory with Democritus. What the Atomists did, provides another example of the method that has just been discussed: like many others, they realised the absurdity of the Eleatic theory, as it had provided no intuitive argument to refute Heraclitus' principle; therefore, they attacked its basic principle, saying that because change and motion surely must be Real, what-is-not \textit{is}, and that what-is moves through it, changing as it does so---i.e., that reality, in their theory, consisted of physical corporeal matter moving about in a physical incorporeal void.

Although none of their original works now remain, the theory they developed, as well as their mode of reasoning, were well-documented through the works of Aristotle; so we are able to present their theory through such quotations. Continuing, then, in line with this brief introduction to the theory, we have, from Simplicius of Cilicia's sixth century \textsc{a.d.} commentary on Aristotle's \textit{Physics}, 28.4-27 \cite{Taylor1999},
\begin{quotation}
Leucippus of Elea or Miletus (both are reported) associated in philosophy with Parmenides, but did not follow the same path as Parmenides and Xenophanes about what there is, but, it seems, the opposite. For they made the universe consist of one changeless, ungenerated, and bounded thing and did not admit that one could even think of what is not, while he posited the atoms, an infinite number of elements in continual motion, and held that they have an infinite number of shapes, since there is no more reason for them to be one shape than another, and that coming to be and change are unceasing among the things that there are. Further, that which is exists no more than that which is not, and both are alike causes of the things that come to be. For positing the nature of the atoms as a solid and a plenum he said that it is what is and that it travels about in the void, which he called `what is not' and said that it is no less than what is.

Similarly his associate Democritus of Abdera posited as principles the plenum and the void, calling the one `what is' and the other `what is not.' They posit the atoms as matter of the things that there are, and generate everything else by their differentiating characteristics. There are three of these, `rhythm,' `turning,' and `contact,' which is to say shape, position and arrangement. For it is natural for like to be affected by like and for things of the same kind to move towards one another and for each of the shapes to be reorganized into a different complex and so make another state. So they claimed that, the principles being infinite, they would readily account for all substances and properties, and what is the cause of anything and how. That is why they say that only those who make their elements infinite can make everything turn out according to their theory. They say that the number of shapes belonging to the atoms is infinite because there is no more reason for them to be one shape than another; for that is the explanation they give of the infinity.
\end{quotation}

Then, from his commentary on Aristotle's \textit{On the Heavens}, 294.33-295.36 (\cite{Aristotle1984}, \textit{F 208 R$^{3}$}), we have the further description,
\begin{quotation}
A few words quoted from Aristotle's \textit{On Democritus} will reveal the line of thought of those men [sc.~the Atomists]:---Democritus thinks the nature of the eternal entities consists of small substances infinite in number; he supposes a place for them, different from them and infinite in extent, and to this he applies the names `void', `nothing', and `the infinite', while to each of the substances he applies the names `thing', `solid', and `existent'. He thinks the substances are so small as to escape our senses, but have all sorts of shapes and figures, and differences of size. From these, then, as from elements, are generated and compounded visible and perceptible masses. The substances are at variance and move in the void because of their dissimilarity as well as the other aforesaid differences, and as they move they collide with each other and interlock in such a way that, while they touch and get close to each other, yet a single substance is never in reality produced from them; for it would be very simple-minded to suppose that two or more things could ever become one. The cause of these substances remaining with one another for some time he ascribes to the bodies fitting into one another and catching hold of one another; for some of them are scalene, others hook-shaped, others concave, others convex, and others have countless other differences. He thinks that they cling to one another and remain together until some stronger necessity arriving from the environment scatters them apart and separates them. He ascribes the genesis and separation opposed to it not only to animals but also to plants, and to worlds, and generally to all perceptible bodies.
\end{quotation}

Finally, we'll consider one of the more elaborate descriptions of the Atomists' theory given by Aristotle himself, which compares it to other contemporary theories, and, incidentally, indicates that the concept of a void was not so much a subsequent realisation of theirs, thereafter understood as a sufficient resolution of the problem posed by the Eleatic theory, as it was actually commonly considered a significant problem throughout Ancient natural philosophy, at least from the time of the Eleatics, and which very much remained a contentious point in Aristotle's time, over a hundred years later---though, probably as a direct result of the Eleatics' reasoning, the Atomists \textit{were} the first to have taken the definitive stance directly opposed to theirs, and to have deduced a theory that was equally consistent with its principles. To be sure: although modern physics provides overwhelming confirmation that vacuous regions of spacetime do actually exist, it is still not a universal belief that spacetime should subsist independent of any material it might contain; so it is, that, according to the principles stated above, we have here taken the same stance that was first truly supported in ancient atomism:---that the incorporeal void of spacetime is to be a genuine component of reality. The relevant quotation, from Aristotle's \textit{On Generation and Corruption}, runs thus \cite{Aristotle1984} (324b25-325b15):
\begin{quotation}
Some philosophers think that the last agent---the agent in the strictest sense---enters in through certain pores, and so the patient suffers action. It is in this way, they assert, that we see and hear and exercise all our other senses. Moreover, according to them, things are seen through air and water and other transparent bodies, because such bodies possess pores, invisible indeed owing to their minuteness, but close-set and arranged in rows---the more transparent the body, the more so.

Such was the theory which some philosophers (including Empedocles) advanced in regard to the structure of certain bodies. They do not restrict it to the bodies which act and suffer action; but combination too, they say, takes place only between bodies whose pores are in reciprocal symmetry. The most systematic theory, however, and one that applied to all bodies, was advanced by Leucippus and Democritus: and, in maintaining it, they took as their starting-point what naturally comes first. 

For some of the older philosophers\footnote{This paragraph and the next outline the argument of the Eleatics.} thought that what is must of necessity be one and immovable. The void, they argue, is not; but unless there is a void with a separate being of its own, what is cannot be moved---nor again can it be many, since there is nothing to keep things apart. And they hold that the view that the universe is not continuous but consists of separate things in contact is no different from the view that there are many (and not one) and a void. For if it is divisible through and through, there is no one, and no many either, but the Whole is void; while to maintain that it is divisible at some points, but not at others, looks like an arbitrary fiction. For up to what limit is it divisible? And for what reason is part of the Whole indivisible, i.e. a \textit{plenum}, and part divided? Further, they maintain, it is equally necessary to deny the existence of motion.

Arguing in this way, therefore, they were led to transcend sense-perception, and to disregard it on the ground that one ought to follow reason; and so they assert that the universe is one and immovable. Some of them add that it is infinite, since the limit (if it had one) would be a limit against the void.\footnote{This was a contribution of Melissus'; see, e.g., \cite{Barnes1982} for further description.}

There were, then, certain thinkers who, for the reasons we have stated, enunciated views of this kind about the truth \ldots Moreover,\symbolfootnote[1]{`One or more arguments against the Eleatic theory appear to have dropped out' (Joachim) [i.e., translator's note].} although these opinions appear to follow logically, yet to believe them seems next door to madness when one considers the facts. For indeed no lunatic seems to be so far out of his senses as to suppose that fire and ice are one: it is only between what \textit{is} right, and what \textit{seems} right from habit, that some people are mad enough to see no difference.

Leucippus, however, thought he had a theory which harmonized with sense-perception and would not abolish either coming-to-be and passing-away or motion and the multiplicity of things. Making these concessions to the phenomena and conceding to the Monists that there could be no motion without a void, he states that void is not-being, and no part of what is is not-being; for what is in the strict sense of the term is an absolute \textit{plenum}. This \textit{plenum}, however, is not one: on the contrary, it is a many infinite in number and invisible owing to the minuteness of their bulk. The many move in the void (for there is a void); and by coming together they produce coming-to-be, while by separating they produce passing-away. Moreover, they act and suffer action wherever they chance to be in contact (for they are not thereby one), and they generate by being put together and becoming intertwined. From the genuinely one, on the other hand, there never could have come-to-be a multiplicity, nor from the genuinely many a one: that is impossible. But just as Empedocles and some of the other philosophers say that things suffer action through their pores, so all alteration and all passion take place in this way, breaking-up (i.e.~passing-away) being affected by means of the void, and so too growth---solids creeping in to fill the void places.

Empedocles too is practically bound to adopt the same theory as Leucippus. For he must say that there are certain solids which, however, are indivisible---unless there are continuous pores all through the body. But this is impossible; for \textit{then} there will be nothing solid beside the pores but all of it will be void. It is necessary, therefore, for his contiguous things to be indivisible, while the intervals between them---which he calls pores---must be void. But this is precisely Leucippus's theory of action and passion.

Such, approximately, are the accounts of the manner in which some things act while others suffer action. And as regards the Atomists, it is not only clear what their explanation is: it is also obvious that it stands in tolerable consistency with the assumptions they employ.
\end{quotation}
He then moves on to consider the inconsistency in Empedocles' theory, before coming back to compare the atomic theory to that of Plato, upon which he concludes, `Thus the comings-to-be and the dissociations result from the indivisibles \textit{according to Leucippus} through the void and through contact (for it is at the point of contact that each of the composite bodies is divisible), but \textit{according to Plato} in virtue of contact alone, since he denies there is a void.'

Therefore, according to the atomic theory, the void would allow a place for every world (according to similar accounts, the theory held that there should have been infinitely many of them in the void, continually coming-to-be and passing-away) to have been in the past, and to move to in the future, while the atoms' continued motion within each world would allow for the change, locomotion, alteration, etc., of things, as we observe. The atoms were required to have been moving through the void since eternity, in order to naturally produce coming-to-be and passing-away, of everything from whole worlds to the smaller bodies within worlds; but this seems to have come at some price, as it suggests that no atom, or aggregate body, should have the ability to actually choose its path within the world, because any change that would thus occur must come about by necessity, and would be as strictly determined as in the Eleatic theory; for so the corporeal past-and-future of the Eleatic theory, having been replaced by a real void, and therefore not physically real itself, was nevertheless only changed to be an equally absolute metaphysical conception. 

However, it is uncertain whether the Atomists did accept such a requirement of strict determinism in their theory, or whether they had in fact allowed the possibility that animate beings, once formed, would be capable of exercising free will, choosing their paths through the void according to reason, as much as by necessity---and arguments to support either side can easily be made from the existing fragments and testimonia; see, e.g., \cite{Taylor1999} for further discussion on this point, which I'll leave off with after two further quotations: first, Diogenes La\"{e}rtius relates Democritus' opinions as having been that \cite{Diogenes1931}
\begin{quotation}
\ldots the atoms are unlimited in size and number, and they are borne along in the whole universe in a vortex, and thereby generate all composite things---fire, water, air, earth; for even these are conglomerations of given atoms. And it is because of their solidity that these atoms are impassive and unalterable. The sun and the moon have been composed of such smooth and spherical masses [\textit{i.e.}~atoms], and so also the soul, which is identical with reason. We see by virtue of the impact of images upon our eyes.

All things happen by virtue of necessity, the vortex being the cause of the creation of all things, and this he calls necessity. The end of action is tranquillity, which is not identical with pleasure, as some by a false interpretation have understood, but a state in which the soul continues calm and strong, undisturbed by any fear or superstition or any other emotion. This he calls well-being and many other names. The qualities of things exist merely by convention; in nature there is nothing but atoms and void space.
\end{quotation}
And secondly, we have the one fragment that is likely attributable to Leucippus: `Nothing happens in vain, but everything from reason and by necessity' \cite{Taylor1999}. The fact that reason, though its existence is by virtue of necessity, as with all else, was nevertheless included in this statement, suggests that it was not stated here redundantly, but actually as a basic cause of happenings.

Compared with Aristotle's later theory of `action' and `passion' to explain interaction within a continuous distribution of matter, i.e.~with no void place anywhere, the original atomic theory, in which individual atoms would be jostled about within an incorporeal void, is far closer to the modern physical conception of reality, as it maintained that things would move about and be seen to do so as the elements of observation would physically traverse space, from one separate body to another. But consider that even to us, who, e.g., accept the concept of molecules of a gas in thermal motion as something natural, this idea in fact remains somewhat foreign; for although we are generally content to think of photons, neutrinos, and other particles travelling through practically empty space, e.g.~\textit{en route} from the Sun to the Earth, we still marvel at the fact that trillions of neutrinos generated by the Sun pass through our bodies every second---indeed, that a proportionally greater number pass through the Earth every second, which impedes them little more than if it were a pure vacuum---\textit{only because they are electrically neutral}, and can therefore only interact with other particles when they pass at inconceivably small distances. To a neutrino, even the Earth is, in effect, a sparse collection of particles separated by huge expanses of void space. 

But consider that even Einstein, through Mach's influence, had kept as close to the Eleatic view as he was able; e.g., in `Relativity and the Problem of Space', he argued \cite{Einstein2001},
\begin{quotation}
\ldots. Mach, in the nineteenth century, was the only one who thought seriously of an elimination of the concept of space, in that he sought to replace it by the notion of the totality of the instantaneous distances between all material points. \ldots

[In Newtonian mechanics], ``physical reality'', thought of as being independent of the subjects experiencing it, was conceived as consisting, at least in principle, of space and time on one hand, and of permanently existing material points, moving with respect to space and time, on the other. The idea of the independent existence of space and time can be expressed drastically in this way: If matter were to disappear, space and time alone would remain behind (as a kind of stage for physical happening).

\ldots 

We are now in a position to see how far the transition to the general theory of relativity modifies the concept of space. In accordance with classical mechanics and according to the special theory of relativity, space (space-time) has an existence independent of matter or field. In order to be able to describe at all that which fills up space and is dependent on the coordinates, space-time or the inertial system with its metrical properties must be thought of at once as existing, for otherwise the description of ``that which fills up space'' would have no meaning.\symbolfootnote[1]{If we consider that which fills space (e.g.~the field) to be removed, there still remains the metric space in accordance with 
\begin{equation}
ds^2={dx_1}^2+{dx_2}^2+{dx_3}^2-{dx_4}^2,\label{M4}
\end{equation}
which would also determine the inertial behaviour of a test body introduced into it.} On the basis of the general theory of relativity, on the other hand, space as opposed to ``what fills space'', which is dependent on the co-ordinates, has no separate existence. Thus a pure gravitational field might have been described in terms of the $g_{ik}$ (as functions of the coordinates), by solution of the gravitational equations. If we imagine the gravitational field, i.e.~the functions $g_{ik}$, to be removed, there does not remain a space of the type [Eq.~(\ref{M4})], but absolutely \textit{nothing}, and also no ``topological space''. For the functions $g_{ik}$ describe not only the field, but at the same time also the topological and metrical structural properties of the manifold. A space of the type [Eq.~(\ref{M4})], judged from the standpoint of the general theory of relativity, is not a space without field, but a special case of the $g_{ik}$ field, for which---for the coordinate system used, which itself has no objective significance---the functions $g_{ik}$ have values that do not depend on the co-ordinates. There is no such thing as an empty space, i.e.~a space without field. Space-time does not claim existence on its own, but only as a structural quality of the field.
\end{quotation}

When we see that even one of the greatest minds of the twentieth century, with all the knowledge of modern physics at his disposal, had propounded a block universe description in which space itself, although it could not be a plenum, should only exist as an accidental property of the more complex, scientifically advanced, four-dimensional field, it should not be difficult to see why, in the fifth century \textsc{b.c.}, aside from the fact that the principles of Atomism gave rise to a comparably more reasonable description of perceptions of Nature, the Atomists' void would have seemed to many to be a gratuitous hypothesis, or, as Einstein still thought of it in his day, an `absurd notion', when there would have been no obvious indication that any such thing as a Void should have independent existence in the world, seen to consist of the Earth, its atmosphere, and all the heavenly bodies. 

If we examine the situation from this perspective, making a further note of consistency, that the Eleatics would have had good reason for believing in the perfect continuity of all matter, which they believed to such an extent that they accepted the consequence that motion and change would have to be denied simply because such things couldn't be Real in the infinite bottleneck, then it seems reasonable that the later abstract metaphysicians, in utilising the tools of their trade, would have seen no good reason that there should have been a void, as, in a very different way, they could sufficiently explain all of the phenomena. For then, even if matter did exist in a real physical place, it could be that this place would be continuously filled with matter, and that change occurred by rearrangement of matter in time, as witnessed by this continually changing arrangement of things. In his \textit{Physics}, Aristotle wrote \cite{Aristotle1984} (208b1-8):
\begin{quotation}
The existence of place is held to be obvious from the fact of mutual replacement. Where water now is, there in turn, when the water has gone out as from a vessel, air is present; and at another time another body occupies this same place. The place is thought to be different from all the bodies which come to be in it and replace one another. What now contains air formerly contained water, so that clearly the place or space into which and out of which they passed was something different from both.
\end{quotation}
Accordingly, e.g., even when a dry sponge is compressed, it would not be compressed into place devoid of matter, but air would have filled parts of its interior, which would then be expressed into the surroundings. Such are the reasons why Aristotle could, one hundred years later, following the Athenian progress in metaphysics, refute the existence of void, as empty place, and propound a theory in which place would be everywhere full, but still physically changing in time. Of course, his theory had its own particular set of problems, such as a requirement to account for the sensation of things that happen at a distance, which he resolved with his theory of action and passion,---as he did with each of the other problems that his theory encountered, again with tolerable consistency, even though that theory is now understood to have been completely false.
\section{On Time\label{sec_OT}}
It is now common knowledge that the corporeal existence supposed in Aristotle's physical theory of Nature was incorrect; and moreover, that the Atomists had been basically on the right track. As I read the accounts of their theory, it seems to me that it is primarily lacking in the way it treats time when accounting for the processes that must take place \textit{through} time. According to certain accounts---e.g., the cosmogony of Leucippus that was summarised by Diogenes La\"{e}rtius \cite{Diogenes1931}---it should be understood that the theory described an infinite three-dimensional void containing infinitely many bounded `worlds', of which our observable Universe would be one. However, according to the modern understanding of the large-scale Universe, due to the past century's cosmological observations, the realistic sense that It must be unbounded suggests that the Atomists' theory should have been more accurate if time had instead been described by a real fourth dimension of the void, through which each `world', unbounded in three dimensions, would move uniformly, always remaining a smooth hypersurface. Furthermore, regarding the consistent physical progression of time, as it must take place in the three-dimensional atomic cosmology, which is formally no different a problem than it was in his own theory, Aristotle recognised that \cite{Aristotle1984} (252a32-b6)
\begin{quote}
\ldots it is a wrong assumption to suppose universally that we have an adequate first principle in virtue of the fact that something always is so or always happens so. Thus Democritus reduces the causes that explain nature to the fact that things happened in the past in the same way as they happen now; but he does not think fit to seek for a principle to explain this `always': so, while his theory is right in so far as it is applied to certain individual cases, he is wrong in making it of universal application. Thus, a triangle always has its angles equal to two right angles, but there is nevertheless an ulterior cause of the eternity, whereas principles are external and have no ulterior cause. Let this conclude what we have to say in support of our contention that there never was a time when there was not motion, and never will be a time when there will not be motion.
\end{quote}
Aristotle's argument for a physical time in which the type of motion he means here---an abstract metaphysical motion that is (to briefly paraphrase the principal conclusions he makes in the final two books of his \textit{Physics} \cite{Aristotle1984}, without including any of the logical argument) numerically one and the same, and simple; that is primary and the measure of all other motions: the continuous rotary motion imparted by a prime mover who alone remains eternally unmoved---is really not relevant for us because he denied the subsistence of the void in principle, having concluded an argument against it sometime earlier on. Instead, we will eventually come to a different argument for physical time that is more relevant to our purposes. For now, it should be sufficient to note that time must be physical, and consider some consequences that this has for the `four-dimensional atomic theory' just described.

It is true that one problem a theory of four real dimensions would have faced, which is possibly also a reason why it was not taken up by the Atomists, is one that arises when extra dimensions need to be added in the derivation of a physical description of Nature, which was recognised by Zeno of Elea, who wrote a collection of paradoxes intended to pose difficulties for any theory that might come against the Eleatic one; the particular one of relevance here, is \cite{Aristotle1984} (209a25), `if everything that exists has a place, place too will have a place, and so on \textit{ad infinitum}.'

However, if the three dimensions of space should be unbounded, and if time should be recognisable as the relative positioning of all matter as the Universe really moves through a physical time-dimension, due to a prior physical property of reality according to which time should continually pass with consistent measure, then this problem can be addressed by considering the only two possibilities that then exist: either time should be the result of a more abstract metaphysical cause, as both the Atomists and Aristotle appear to have thought it was, witnessed through the continual rearrangement of the Universe (whether or not there would be physical void separating atoms in space), which still needs to be addressed in the physical theory, as Aristotle seems to be the first to have recognised; or, time is an extra dimension of a physical void, through which every body in the Universe moves continuously, as they also may rearrange themselves in space (which, again, need not be full). But these two cases are obviously not formally distinct, as they both describe time as the motion of a universe through something physical and incorporeal, and the choice, whether to conceive of that as abstract or not, is purely arbitrary.

However, there is another, accidental sense of time in this one formal description: the non-real (purely \textit{ideal}) one that the Atomists probably had in mind, which describes the past, present, and future events that happen in the Universe, as they actually were, are, and will be. With the addition of this dimension, the physical description of all events that would occur as a three-dimensional universe really moved through a real four-dimensional void, would not be four-dimensional (as they would be if that real fourth dimension were neglected, and only the events occurring in the universe through time were taken into account), but actually five-dimensional; however, because the Universe is still supposed to be three-dimensional, the set of events that occur in the Universe itself, as it moves through that fourth real dimension, would only be a four-dimensional slice of that five-dimensional description of everything that ever was and will be.

Actually, the fact that this ideal time dimension exists in the description of any such theory---i.e., even in Aristotle's theory, in which time is the uniform motion imparted on everything by a prime mover, which takes place through an abstract metaphysical dimension---implies a second option, which was left out above by requiring \textit{real} movement through the fourth physical dimension; viz., that real motion does not actually need to take place, but could really be that same accidental dimension in which all past and future events in the Universe occur---the physical description of all occurrences in the Universe---which is simply promoted in status, to be the actual fourth dimension of physical reality---as it was in the Eleatic theory, and as it is described in the original interpretation of relativity theory: a physical four-dimensional block.

So we have two formally distinct descriptions of Time, according to the reasonings of the ancient Greek philosophers. However, as shown in the detailed discussion in \S\S~\ref{sec_OJGKD} -- \ref{sec_ARP}, relativity theory cannot possibly distinguish between these two interpretations of physical reality, because the four-dimensional continuum of events that the formal theory is used to describe---viz., \textit{spacetime}---can always be thought of, either as a block, like the Eleatic One, or as the set of events occurring on a well-defined three-dimensional space dynamically moving through a four-dimensional void, according to Einstein's equation, as in our `four-dimensional atomic theory'; then, it is an immediate consequence that in the latter case the fifth (physical) dimension is not required to be \textit{real}, as the regular evolution of all four dimensions of reality must progress according to the formal dynamics that need to be assumed in the physical theory, and Zeno's infinite sequence of real places is no more required than it is in the Eleatic theory.

Now, it is important to investigate in more detail the physical reasons why Time actually needs to be a physical dimension, as It is formally described by the theories of, e.g., Aristotle, Newton, and Einstein; viz., as the critical philosophical argument that Time must be somehow physical---i.e., subject to the fundamental Laws of physics,---originates in the fact, common to both the pre-relativistic and relativistic time concepts, that Its passage is always \textit{perceived} with consistent measure. A sufficient argument for this was already given---though he did not draw the same natural conclusion from his reasonings as we shall---by Augustine of Hippo, near the end of the fourth century \textsc{a.d.}, in the eleventh Book of his \textit{Confessions}, where he wrote of his personal reflections on the Nature of Time, which Russell described as `a very admirable relativistic theory of time' \cite{Russell1946}. Augustine begins with a speculation on the nature of God's own existence, which Russell paraphrases \cite{Russell1946}:---
\begin{quotation}
Why was the world not created sooner? Because there was not ``sooner.'' Time was created when the world was created. God is eternal, in the sense of being timeless; in God there is no before and after, but only an eternal present. God's eternity is exempt from the relation of time; all time is present to Him at once. He did not \textit{precede} His own creation of time, for that would imply that he was in time, whereas He stands eternally outside the stream of time.
\end{quotation}
This, Russell points out, is what leads Augustine to his `admirable relativistic theory':
\begin{quotation}
``What, then, is time?'' he asks. ``If no one asks of me, I know; if I wish to explain to him who asks, I know not.'' Various difficulties perplex him. Neither past nor future, he says, but only the present, really \textit{is}; the present is only a moment, and time can only be measured while it is passing. Nevertheless, there really is time past and future. We seem to be led into contradictions. The only way Augustine can find to avoid these contradictions is to say that past and future can only be thought of as present: ``past'' must be identified with memory, and ``future'' with expectation, memory and expectation being both present facts. There are, he says, three times: ``a present of things past, a present of things present, and a present of things future.'' ``The present of things past is memory; the present of things present is sight; and the present of things future is expectation.'' To say that there are three times, past, present, and future is a loose way of speaking.

He realizes that he has not really solved all difficulties by this theory. \ldots But the gist of the solution he suggests is that time is subjective: time is in the human mind, which expects, considers, and remembers. It follows that there can be no time without a created being, and that to speak of time before the Creation is meaningless.

I do not myself agree with this theory, in so far as it makes time something mental. But it is clearly a very able theory, deserving to be seriously considered. I should go further and say that it is a great advance on anything to be found on the subject in Greek philosophy. It contains a better and clearer statement than Kant's of the subjective theory of time---a theory which, since Kant, has been widely accepted among philosophers.
\end{quotation}
Russell goes on to conclude that, since Augustine's theory anticipated not only Kant's theory of time, but Descartes' \textit{cogito}, he deserves a high place as a philosopher. But what Russell fails to relate here, which must be of far greater interest to the physicist, is the importance of the evidences Augustine presents, which lead him to his conclusion. For everyone knows that any conclusion reasoned inductively is not necessarily the only conclusion that \textit{can} be drawn from such a process; and, as Maxwell wrote, `It is of great advantage to the student of any subject to read the original memoirs on that subject, for science is always most completely assimilated when it is in the nascent state' \cite{Maxwell1892}. In this case, it should actually become clear to any relativist, that in his loose way of speaking about past, present, and future, as they relate to the measurements of temporal durations by a single observer, Augustine had anticipated modern physics in providing a good basis for certain conceptions of \textit{spacetime} and \textit{proper time}; therefore, his theory is far more relativistic than Russell leads us to believe, as it appears his comment was probably meant to address, instead, the less interesting fact that, generally speaking, the eternal realm of God through which our minds progress, is essentially the reality to which Einstein had attributed the description given by relativity theory. 

The passages of Augustine's text relating these progressive ideas, are \cite{Augustine1912}:
\begin{quotation}
\ldots we take such measure of the times in their passing by, as we may be able to say, this time is twice so much as that one; or, this is just so much as that: and so of any other parts of time, which be measurable. We do therefore, as I said, take measure of the times as they are passing by. And if any man should now ask me: How knowest thou? I might answer, I do know, because we do measure them: for we cannot measure things that are not; and verily, things past and to come are not. But for the present time now, how do we measure that, seeing it hath no space? We measure it therefore, even whilst it passeth, but when it is past, then we measure it not: for there will be nothing to be measured. But from what place, and by which way, and whitherto passes this time while it is a measuring? Whence, but from the future? Which way, but through the present? Whither, but into the past? From that therefore, which is not yet: by that, which hath no space: into that, which is not still. Yet what is it we measure, if not time in some space? For we use not to say, single, and double, and triple, and equal, or any other way we speak of time, but with references still to the spaces of times. In what space therefore do we measure the time present? Whether in the future space, whence it is passing? Or in the present, by which it is passing? But no-space we do not measure. Or in the past, to which it passeth? But neither do we measure that which is not still.

\ldots

\ldots. I will not therefore demand now what that should be which is called day: but, what time should be, by which we measuring the circuit of the sun, should say, that he had then finished it in half the time he was wont to do, if so be he had gone it over in so small a space as twelve hours come to: and when upon comparing of both times together, we should say, that this is but a single time, and that a double time, notwithstanding that the sun should run his round from east to east sometimes in that single time, and sometimes in that double time. \ldots I perceive time therefore to be a certain stretching. But do I perceive it, or do I seem to perceive it? \ldots

\ldots. Seeing therefore the motion of a body is one thing, and that by which we measure how long it is, another thing; who cannot now judge which of the two is rather to be called time? For and if a body be sometimes moved uncertainly, and stands still other sometimes; then do we measure, not his motions only, but his standing still too: and we say, It stood still as much as it moved; or It stood still twice or thrice so long as it moved; or any other space which our measuring hath either perfectly taken, or guessed at, more or less, as we use to say. Time therefore is not the motion of a body.

\ldots

\ldots it seems to me, that time is nothing else but a stretching out in length; but of what, I know not, and I marvel, if it be not of the very mind. For what is it, I beseech thee, O my God, that I now measure, whereas I say, either at large, that this is a longer time than that: or, more particularly, that this is double to that? I know it to be time that I measure: and yet do I neither measure the time to come, for that is not yet: nor time present, because that is not stretched out in any space: nor time past, because that is not still. Is it the times as they are passing, not as they are past? For so I was saying.
\end{quotation}

From these speculations, the progression of Augustine's thoughts continue towards that which must have seemed to him to be the simplest conclusion, and therefore, according to the need for logical economy, the most acceptable one: viz., that the `space of time' might be mental. We are quite sure now, however, that this is not the case,---that time itself is also a real physical dimension of a greater subsistent field; but it is also obvious that a physical reality of time is the natural conclusion of Augustine's speculations, regardless of whether that would be only mental. 

On the other hand, his argument in no way implies that he has thought of any difference in observation of the passages of time measured by two relatively moving observers; nor that time and space should form something we could recognise as a topologically connected, four-dimensional space; thus, Augustine's reasoning applies sufficiently well to the concept in Newton's theory, that \cite{Newton1966}
\begin{quote}
Absolute, true, and mathematical time, of itself, and from its own nature, flows equably without relation to anything external, and by another name is called duration: relative, apparent, and common time, is some sensible and external (whether accurate or unequable) measure of duration by the means of motion, which is commonly used instead of true time; such as an hour, a day, a month, a year.
\end{quote} 

This should, however, not be thought of as a mark against the theory, but merely as an incompleteness, as his idea is in fact more basic---merely that there is, all ways, a prior, consistently measurable duration---which is why it pertains indiscriminately to both the pre-relativistic and relativistic time concepts; i.e., we can say that time, according to Augustine's theory, is something that \textit{is} extended, in whatever capacity that may be; and that it is measurable by any observer, who, for themself, must be sensible of a dynamic present, as well as its associated past and future, in which durations of processes within proper three-dimensional space can be observed, but which exists prior to any measurable occurence, underlying the very idea of process---whether that would be observed as active or inertial.

Thus, by the very nature of his argument, in which he considers time only from the subjective standpoint of personal observation, and does not consider how its measurement would differ according to the perspective of an external body---e.g., which might move some distance, then sit awhile, then perhaps spin in place,---regardless of the fact that he did not realise the additional complication of relative motion, Augustine did find a good description of time, witnessed from the proper perspective as some basic physically extended aspect of Nature, which may or may not be absolute, and thus could, as well as not, be used to form a basis of an underlying four-dimensional, connected spacetime. Augustine's contribution to the natural philosophy of time can therefore be stated: He solidified Heraclitus' principle of uniform flux according to a further rationalisation of common sense which led him to realise that this natural progression through time always appears to have consistent measure, and must therefore be physical, laid out along a physical dimension.

So we see the relevant issue; viz., that Augustine's theory, though generally applicable to more developed physical theories, actually did not go far enough to be considered a completed theory of time. And, of course, we know that the final essential advancement that was made on this singular perspective came only fifteen hundred years later, after the relativity of inertia had emerged from Renaissance thought, and, furthermore, following the development of a theory of light which culminated in the precision measurement, against conventional reason, that its speed should be the same constant and finite value in all frames,---when Einstein, in postulating natural principles relating to the kinematics of these observations, realised the physical consequences pertaining to measures of time and simultaneity, when two relatively moving observers' prespectives would then be considered in relation to one another; viz., that the two observers would not agree on the simultaneity of events perceived in their proper frames. 
\section{The Original Interpretation of Spacetime\label{sec_OIS}}
The relativity of simultaneity requires that the proper coordinate system of every relatively moving frame must provide a relatively unique description of reality; and the differentiating factor between the interpretation of relativity theory to be discussed in this section, and the new one that is described in Chapter~\ref{chap_DIRT}, is in fact how that reality should be formally identified in the mathematical description, if it should actually be dynamic.

The original interpretation of spacetime, which has remained a central aspect of the theory, was first described in 1908, by Einstein's former mathematics teacher from his days in Z\"{u}rich, Hermann Minkowski, who, at a conference in Cologne, proposed both: that relativity theory should be taken to mean that space and time are not disjoint, but actually form a four-dimensional continuum; and, subsequently---in a metaphysical statement equvalent to one for which McTaggart argued at length in `The Unreality of Time', published the following month \cite{McTaggart1908},---that this should be the True aspect of Reality, in which all events exist immediately, throughout the entire continuum, once and for all. In his words \cite{Einstein1952},
\begin{quotation}
\ldots space by itself, and time by itself, are [henceforth] doomed to fade away into mere shadows, and only a kind of union of the two will preserve an independent reality.

\ldots

\ldots. A point of space at a point of time, that is, a system of values $x$, $y$, $z$, $t$, I will call a \textit{world-point}. The multiplicity of all thinkable $x$, $y$, $z$, $t$ systems of values we will christen the \textit{world} \ldots. Not to leave a yawning void anywhere, we will imagine that everywhere and everywhen there is something perceptible \ldots. Then we obtain, as an image, so to speak, of the everlasting career of the substantial point, a curve in the world, a \textit{world-line}, the points of which can be referred unequivocally to the parameter $t$ from $-\infty$ to $+\infty$. \ldots

\ldots

\ldots. We should then have in the world no longer \textit{space}, but an infinite number of spaces, analogously as there are in three-dimensional space an infinite number of planes. Three-dimensional geometry becomes a chapter in four-dimensional physics. Now you know why I said at the outset that space and time are to fade away into the shadows, and only a world in itself will subsist. 

\ldots

Lorentz called the $t'$ combination of $x$ and $t$ the local time of the electron in uniform motion, and applied a physical construction of this concept, for the better understanding of the hypothesis of contraction. But the credit of first recognizing clearly that the time of the one electron is just as good as that of the other, that is to say, that $t$ and $t'$ are to be treated identically, belongs to A.~Einstein \cite{Einstein1905a}. Thus time, as a concept unequivocally determined by phenomena, was first deposed from its high seat. Neither Einstein nor Lorentz made any attack on the concept of space, perhaps because in the \ldots special transformation, where the plane of $x'$, $t'$ coincided with the plane of $x$, $t$, \textit{an interpretation is possible} by saying that the $x-$axis of space maintains its position. [Emphasis added.]
\end{quotation} 
Thus, in the tradition of the Eleatics, Minkowski placed the mathematical description of any physical \textit{process}, whether `active' or `inertial', in the same realm that St.~Augustine had placed God. As he appears to have realised roughly the additional possibility of the interpretation that will be argued for in Chapter~\ref{chap_DIRT}, we might wonder whether his reasons for decrying this interpretation should have had something to do with the simple facts, that the interpretation he chose to propound is mathematically more neat, or that he---again, in the Eleatic tradition---clearly abhorred the possibility of a reality in which there could physically be a `yawning void', and so impressed that bias on natural science. Regardless, the interpretation he did decide on is the one that stuck---at least, it did with Einstein.

For it is well-documented that although Einstein did not immediately take to Minkowski's interpretation, he did come around to accepting it, as he realised the necessity of four-dimensional spacetime in order to arrive at the general theory. Pais \cite{Pais1982} notes that Einstein once remarked to V.~Bargmann that he `regarded the transcriptions of his theory into tensor form as ``\"{u}berfl\"{u}ssige Gelehrsamkeit,{''} (superfluous [erudition]).' It seems likely that this sentiment comes from one that Einstein is known to have carried with him from his youth; e.g., in \textit{Albert Einstein: A Documentary Biography} \cite{Seelig1956}, Carl Seelig provides a quotation from Einstein's autobiography, which seems to me to shed some more light on the comment made to Bargmann: 
\begin{quotation}
\ldots my intuition in the mathematical field was not strong enough to be able to distinguish with basic conviction the fundamentally important from the rest of the more or less dispensable erudition. Moreover my interest in acquiring a knowledge of Nature was infinitely stronger and as a student it was not clear to me that the approach to a deeper knowledge of the principles of physics was bound up with the most intricate mathematical methods. This only dawned on me after years of independent scientific work.

\ldots In [physics], however, I soon learned to sense those [data] in particular that could lead to fundamentals and to overlook a host of others that filled the mind and distracted from the essential.
\end{quotation}
Seelig then goes on to describe Minkowski's interpretation, noting that, in 1916, Einstein had paid tribute to its contribution to the development of general relativity theory; after this, however, he notes that `Einstein once joked about the mathematical finesse with which his work was investigated and expanded. He said with a sigh: ``Since the mathematicians have attacked the relativity theory, I myself no longer understand it any more.{''}'\footnote{This redundancy makes the English translation of Einstein's remark seem strikingly poor. Furthermore, there is total ambiguity as to whether Seelig heard this directly from Einstein, or whether this information had come from some other source, which he neglected to reference. Therefore, it is important to refer back to the original german text, in \textit{Albert Einstein und die Schweiz} \cite{Seelig1952}, which reads, `Scherzhaft soll Einstein einmal angesichts der mathematischen Finessen, mit dennen seine Arbeiten untersucht und erweitert wurden, geseufzt haben: {\flqq}Seit die Mathematiker ueber die Relativitaetstheorie hergefallen sind, verstehe ich sie selbst nicht mehr!{\frqq},' which is more accurately translated as, `Once, in light of the mathematical finesse with which his work had been investigated and expanded, Einstein is supposed to have sighed jokingly, ``Since the mathematicians have descended upon the relativity theory, even I no longer understand it.{''}'} 

This is all closely related to something Philipp Frank had written five years earlier, in \textit{Einstein: His Life and Times} \cite{Frank1947}:
\begin{quotation}
Einstein was at first skeptical about the use of very advanced mathematics in developing physical theories. When in 1908 Minkowski showed that Einstein's special theory of relativity could be formulated very simply in the language of four-dimensional geometry, Einstein had regarded this as the introduction of an involved formalism by which it became rather more difficult to grasp the physical content of the theory. When Max von Laue, in the first comprehensive book on Einstein's relativity theory, presented it in a very elegant mathematical form, Einstein remarked at that time jokingly: ``I myself can hardly understand Laue's book.''

The center of German mathematical teaching and research during this period was the University of G\"{o}ttingen. Minkowski taught there, and the mathematical formulation of the relativity theory had begun there. Einstein once remarked playfully: ``The people in G\"{o}ttingen sometimes strike me, not as if they wanted to help one formulate something more clearly, but instead as if they wanted only to show us physicists how much brighter they are than we.'' Nevertheless, the greatest mathematician in G\"{o}ttingen, David Hilbert, realized that while Einstein did not care for the superfluous formal difficulties in mathematics, he did know how to use mathematics where it was indicated. Hilbert once said: ``Every boy in the streets of our mathematical G\"{o}ttingen understands more about four-dimensional geometry than Einstein. Yet, despite that, Einstein did the work and not the mathematicians.''
\end{quotation}

It is clear that Einstein did not immediately intuitively accept the mathematical formalism that he would eventually use to develop the general theory; but regardless of how he came to accept it, we know that he did eventually entirely support this interpretation of physical reality. For example, when he wrote of the innovation of relativity theory in \textit{The Meaning of Relativity}, which was first published in 1921, he said \cite{Einstein1945},
\begin{quotation}
In the pre-relativity physics space and time were separate entities \ldots. One spoke of points of space, as of instants of time, as if they were absolute realities. It was not observed that the true element of the space-time specification was the event specified by the four numbers $x_1$, $x_2$, $x_3$, $t$. The conception of something happening was always that of a four-dimensional continuum; but the recognition of this was obscured by the absolute character of the pre-relativity time. Upon giving up the hypothesis of the absolute character of time, particularly that of simultaneity, the four-dimensionality of the time-space concept was immediately recognized. It is neither the point in space, nor the instant in time, at which something happens that has physical reality, but only the event itself. There is no absolute (independent of the space of reference) relation in space, and no absolute relation in time between two events, but there is an absolute (independent of the space of reference) relation in space and time.
\end{quotation}
Following this description, he speaks of the relativistic field as a `four-dimensional continuum of events'; and according to this specification, everything that has or ever will happen, must exist as one absolute four-dimensional field---that which Minkowski had dubbed, in 1908, the `absolute world'.\footnote{Minkowski even went as far as to say that the word \textit{relativity-postulate} seemed feeble, and that, since it `comes to mean that only the four-dimensional world in space and time is given by phenomena \ldots, I prefer to call it the \textit{postulate of the absolute world} (or briefly, the world-postulate)' \cite{Einstein1952}.} 

The interpretation of reality---i.e., the ontology---Einstein thought was most objectively represented by relativity theory, is more decisively stated in \textit{The Evolution of Physics} \cite{Einstein1938}, published in 1938, which he wrote with Leopold Infeld. There, Minkowski's proposal of the absolute relativistic spacetime is clearly supported over a dynamic one. They begin their discussion by considering two different descriptions of a particle's motion in one dimension, which they claim to be the two possible objective interpretations prior to considering the results of relativity theory:
\begin{quotation}
We remember the picture of the particle changing its position with time in the one-dimensional space. We picture the motion as a sequence of events in the one-dimensional space continuum. We do not mix time and space, using a \textit{dynamic} picture in which positions \textit{change} with time.

But we can picture the same motion in a different way. We can form a \textit{static} picture, considering the curve in the two-dimensional time-space continuum. Now the motion is represented as something which \textit{is}, which exists in the two-dimensional time-space contimuum, and not as something which changes in the one-dimensional space continuum.

Both these pictures are exactly equivalent and preferring one to the other is merely a matter of convention and taste.

Nothing that has been said here about the two pictures of the motion has anything whatever to do with the relativity theory. Both representations can be used with equal right, though classical physics favored rather the dynamical picture describing motion as happenings in space and not as existing in time-space. But the relativity theory changed this view. It was distinctly in favor of the static picture and found in this representation of motion as something existing in time-space a more convenient and more objective picture of reality. We still have to answer the question: why are these two pictures, equivalent from the point of view of classical physics, not equivalent from the point of view of the relativity theory?

The answer will be understood if two co-ordinate systems moving uniformly, relative to each other, are again taken into account.
\end{quotation}
Then, they use the relativity of simultaneity to argue that relativity theory is only consistent with the static picture. In concluding this description of spacetime, however, they write, `Even in the relativity theory we can still use the dynamic picture if we prefer it. But we must remember that this division into time and space has no objective meaning since time is no longer ``absolute.'' We shall still use the ``dynamic'' and not the ``static'' language in the following pages, bearing in mind its limitations.'

In fact, we can use the dynamic description of spacetime in relativity theory because a more objective dynamic description \textit{does} exist, with a clear representation in the theory, which could only have been considered `less convenient' by any modern relativist because it had not been formally addressed, following Minkowski's peremptive defamation. But, in order to see that this possibility clearly must be allowed, every aspect of the problem needs to be better understood---in particular, how the relativity of simultaneity is described in that picture. But we can see already, that Einstein and Infeld's description is founded on a prior loss of the objectivity they want to claim, as they allow only the two possible descriptions, of a sequence of events in one dimension, or prior determination of a static two-dimensional spacetime continuum of events, and explicitly deny the possibility of the relative motion of two particles through a one-dimensional universe, as it moves through another dimension.

But Einstein either never did see how this could work---how a truly dynamic picture could be reconciled with the theory, as I've begun to explain---or else he, too, did not like the implications of this for the nature of reality. Apparently, he had really convinced himself of the truth of the strictly determined block universe, and held fast to that belief, to the same extent that the Eleatics had; for in Karl Popper's autobiography, he wrote that he had discussed indeterminism at length with Einstein, and `tried to persuade him to give up his determinism, which ammounted to the view that the world was a four-dimensional block universe in which change was a human illusion, or very nearly so. (He agreed that this had been his view, and while discussing it I called him ``Parmenides''.)' \cite{Popper2002}; and, in the fifth appendix to \textit{Relativity: The Special and General Theory}, which he added in 1952, he wrote \cite{Einstein2001},
\begin{quotation}
According to the special theory of relativity \ldots the sum total of events which are simultaneous with a selected event exist, it is true, in relation to a particular inertial system, but no longer independently of the choice of inertial system. The four-dimensional continuum is now no longer resolvable objectively into sections, which contain all simultaneous events; ``now'' loses for the spatially extended world its objective meaning. It is because of this that space and time must be regarded as a four-dimensional continuum that is objectively unresolvable, if it is desired to express the purport of objective relations without unnecessary conventional arbitrariness.

\ldots. The four-dimensional structure (Minkowski-space) is thought of as being the carrier of matter and of the field. Inertial spaces with their associated times are only privileged four-dimensional coordinate systems that are linked together by the linear Lorentz transformations. Since there exist in this four-dimensional structure no longer any sections which represent ``now'' objectively, the concepts of happening and becoming are indeed not completely suspended, but yet complicated. It appears therefore more natural to think of physical reality as a four-dimensional existence, instead of, as hitherto, the \textit{evolution} of a three-dimensional existence.

This rigid four-dimensional space of the special theory of relativity \ldots
\end{quotation}
And so finally, only weeks before Einstein's death, in a letter commemorating his friend, Michel Besso, he wrote, `Now he has preceded me a little by parting from this strange world. This means nothing. To us believing physicists the distinction between past, present, and future has only the significance of a stubborn illusion' \cite{Foelsing1997}.

But while Einstein had committed the Eleatic error of not critically investigating the legitimacy of his physical principles in order to reconcile his theory with Reality, and accepted the formal implication of something which is not merely \textit{unwelcome}, but is actually \textit{totally preposterous}, the more common practice amongst theoretical physicists has been to take the path of disbelieving, or even entirely ignorant, sceptics, seemingly unable even to fathom the fact that the formalism we owe to Minkowski does---and was actually intended to!---describe a strictly determined block eternity, by simply ignoring that formal requirement and using the theory instead to describe paradoxical physical processes subjectively, as they occur simultaneously in any given frame.

An interesting scientific speculation which also relates to the Eleatic worldview, which was published less than a century before relativity theory, is the strict determinist theory stated in Pierre-Simon Laplace's \textit{Essai philosophique sur les Probabilit\'{e}s}, which was published as an introduction to his great \textit{Th\'{e}orie analytique des Probabilit\'{e}s} in later editions. There, after a brief introduction, he opens the section, `De la Probabilit\'{e}', with the statement \cite{Laplace1951}, 
\begin{quotation}
All events, even those which on account of their insignificance do not seem to follow the great laws of nature, are a result of it just as necessarily as the revolutions of the sun. In ignorance of the ties which unite such events to the entire system of the universe, they have been made to depend upon final causes or upon hazard, according as they occur and are repeated with regularity, or appear with regard to order; but these imaginary causes have gradually receded with the widening bounds of knowledge and disappear entirely before sound philosophy, which sees in them only the expression of our ignorance of the true causes.

Present events are connected with preceding ones by a tie based upon the evident principle that a thing cannot occur without a cause which produces it. This axiom, known by the name of \textit{the principle of sufficient reason}, extends even to actions which are considered indifferent; the freest will is unable without a determinative motive to give them birth; if we assume two positions with exactly similar circumstances and find that the will is active in the one and inactive in the other, we say that its choice is an effect without a cause. It is then, says Leibniz, the blind chance of the Epicureans. The contrary opinion is an illusion of the mind, which, losing sight of the evasive reasons of the choice of the will in indifferent things, believes that choice is determined of itself without motives.

We ought then to regard the present state of the universe as the effect of its anterior state and as the cause of the one which is to follow. Given for one instant an intelligence which could comprehend all the forces by which nature is animated and the respective situation of the beings who compose it---an intelligence sufficiently vast to submit these data to analysis---it would embrace in the same formula the movements of the greatest bodies of the universe and those of the lightest atom; for it, nothing would be uncertain and the future, as the past, would be present to its eyes. The human mind offers, in the perfection which it has been able to give to astronomy, a feeble idea of this intelligence. Its discoveries in mechanics and geometry, added to that of universal gravity, have enabled it to comprehend in the same analytical expressions the past and future states of the system of the world. Applying the same method to some other objects of its knowledge, it has succeeded in referring to general laws observed phenomena and in foreseeing those which given curcumstances ought to produce. All these efforts in the search for truth tend to lead it back continually to the vast intelligence which we have just mentioned, but from which it will always remain infinitely removed.
\end{quotation}

The intelligence described here is now commonly called \textit{Laplace's demon}. It is interesting to note the particular stance he has assumed regarding spacetime theory: viz., he must have thought of time as it is described in the Newtonian theory, as a disjoint, abstract dimension of physical reality; it is also clear that he has not, as Minkowski eventually did, made a final commitment to conceive of the events that would take place in time as really existing in a pre-determined state, as an eternal block spacetime; rather, he probably thought of the time-dimension, if he ever thought of it as more than a theoretical construct, as an other metaphysical dimension used to describe the events that did occur, are presently occurring, and will eventually occur, in the mode of the ancient Atomists and Aristotle, along with its formal definition in Newtonian physics. 

However, given his obvious belief in the truth of an ideal God-like intelligence that could conceive of, and formulate, all of eternity, as it would necessarily be laid out when acts of free will would be relegated to the realm of infinitely complex probabilistic events, the ontological distinction between either spacetime conception becomes, in effect, what Einstein should have referred to as `superfluous erudition'. In this way, Laplace falls into the Eleatic tradition, as a precursor of Minkowski.

But a last remark, which relates back to our discussion in \S~\ref{sec_ANE}, is warranted by this over-zealous advocacy of any scientific evidence for absolute determinism: it is reasonable to assume that every physical process in Reality must obey Natural Laws which apply objectively throughout space and time, which has been done here as well; but it is nevertheless too bold an inductive leap, to suppose, from the observation of this natural order in various phenomena, that nothing should be left up to free will, which common sense tells us we must possess. Therefore, in the works of Laplace and Minkowski alike, we should recognise an excessive appeal to logical economy in discerning a mathematical description of Nature, which led each of them to propose narrow-minded theories with an unacceptable loss of objectivity, as both had, for this reason and none other, denied basic truths founded on common sense.